\documentclass[twocolumn,10pt,a4paper]{article}
\usepackage[margin=2cm, columnsep=0.6cm]{geometry}
\usepackage{amsmath,amssymb,amsthm}
\usepackage{mathtools}
\usepackage{bm}
\usepackage{graphicx}
\usepackage{hyperref}
\usepackage{natbib}
\usepackage{physics}
\usepackage{booktabs}
\usepackage{array}
\usepackage{multirow}
\usepackage{enumitem}
\usepackage{xcolor}
\usepackage{float}
\usepackage{algorithm}
\usepackage{algpseudocode}
\usepackage[utf8]{inputenc}
\usepackage[T1]{fontenc}
\usepackage{lmodern}
\usepackage{microtype}

%% Custom commands
\newcommand{\kB}{k_{\mathrm{B}}}
\newcommand{\Hcat}{\mathcal{H}}
\newcommand{\dcat}{d_{\mathrm{cat}}}
\newcommand{\taup}{\tau_p}
\newcommand{\Gcond}{\mathcal{G}}
\newcommand{\Back}{\mathcal{B}}
\newcommand{\Muij}{\mu_{ij}}

%% Theorem environments
\newtheorem{theorem}{Theorem}[section]
\newtheorem{lemma}[theorem]{Lemma}
\newtheorem{proposition}[theorem]{Proposition}
\newtheorem{corollary}[theorem]{Corollary}
\newtheorem{definition}[theorem]{Definition}
\newtheorem{axiom}[theorem]{Axiom}
\newtheorem{remark}[theorem]{Remark}
\newtheorem{example}[theorem]{Example}

\title{\textbf{A Unified Cellular Circuit Model:\\[4pt]
Fuzzy State Representation, Backward Trajectories,\\
and Categorical State Propagation}}

\author{Kundai Farai Sachikonye\\[4pt]
\small Technical University of Munich, School of Life Sciences\\
\small \texttt{kundai.sachikonye@wzw.tum.de}}

\date{\today}

\begin{document}
\maketitle

%%=============================================================================
\begin{abstract}
%%=============================================================================
We present a unified framework modelling the biochemical interior of a living
cell as an electrical circuit graph, deriving all structural analogies
rigorously from first principles of thermodynamics, information theory, and
stochastic kinetics.  The central epistemological observation motivating the
framework is \emph{epistemic blindness}: a cell possesses no internal mechanism
capable of distinguishing its own healthy state from a diseased state, because
all intracellular information processing consists of forward reaction dynamics
that faithfully propagate whatever molecular configuration is currently present.
A faithful simulation of those dynamics therefore adds zero new information
relative to the cell itself; if simulation could reveal disease, cells would
self-diagnose---which they cannot.  This constraint motivates a shift from
forward simulation to \emph{trajectory completion}: determining the complete
state of every molecular node from partial observations by exploiting the
topological and thermodynamic structure of the reaction network.

In the circuit model each molecular species is a node carrying a potential
equal to its chemical potential, interpreted information-theoretically as
categorical depth $\Hcat_i = -\log_2 P(\sigma_i)$.  Reactions are edges with
conductance proportional to the forward rate constant times substrate
concentration.  Kirchhoff's current law (KCL) and voltage law (KVL) analogs
are derived exactly from stoichiometric mass balance and thermodynamic cycle
consistency (detailed balance).  Node states are initialised as fuzzy
membership functions over concentration space, and the Zadeh extension
principle lifts the Kirchhoff equations to fuzzy arithmetic.  The trajectory
completion algorithm iterates three constraint classes---Kirchhoff propagation,
maximum-a-posteriori backward-trajectory inference via the Viterbi algorithm,
and thermodynamic feasibility---and is shown to converge to a unique crisp
state assignment by the Banach fixed-point theorem on the space of fuzzy state
tuples equipped with the Hausdorff metric.

We prove that the MAP backward trajectory of any node is \emph{time-invariant}
for autonomous kinetics: it depends on the current state and network topology
but not on the absolute time at which observation occurs.  This trajectory is
therefore a \emph{categorical address}---a fixed geometric identity of the node
within the network's state space.  Categorical state information propagates
through the network at signal velocity $v_s \sim \sqrt{\kappa/m_{\mathrm{eff}}}$,
exceeding Fickian drift by three to six orders of magnitude, rendering the
network informationally well-mixed on timescales shorter than any metabolic
reaction.  Applications to cell-state instantiation from partial omics data,
disease detection via trajectory escape from the healthy attractor basin, and
rational drug design as targeted edge-conductance modification are developed in
detail.  The framework unifies flux balance analysis, metabolic control
analysis, and stochastic kinetics within a single information-geometric picture.
\end{abstract}

%%=============================================================================
\section{Introduction}
\label{sec:intro}
%%=============================================================================

\subsection{The Simulation Problem}
\label{sec:intro-sim}

The dominant paradigm in quantitative cell biology is \emph{forward numerical
simulation}.  Given an initial condition $\mathbf{x}_0\in\mathbb{R}^N$
describing the concentrations of $N$ molecular species, and a system of
ordinary (or stochastic) differential equations
\begin{equation}
  \dot{\mathbf{x}} = \mathbf{S}\,\mathbf{v}(\mathbf{x}),
  \label{eq:ode}
\end{equation}
where $\mathbf{S}$ is the $N\times R$ stoichiometric matrix and
$\mathbf{v}:\mathbb{R}^N\to\mathbb{R}^R$ is the vector of reaction rates, one
integrates forward in time to obtain trajectories
$\mathbf{x}(t)$~\cite{Murray2002,Klipp2005}.  This paradigm has produced
landmark results in understanding biochemical oscillations, bistability, and
signal transduction~\cite{Tyson2003,Strogatz1994}.

Yet it contains a deep epistemic limitation.  Consider what forward simulation
actually computes: given the state at time $t$ it produces the state at
$t+\Delta t$.  The simulation is, by construction, a faithful representation
of the system's own dynamics.  It does not transcend those dynamics; it embodies
them.  Consequently, a faithful simulation of a cell is informationally
equivalent to the cell itself.

This observation has a precise information-theoretic formulation.  Let $\Omega$
denote the set of all trajectories $\mathbf{x}(\cdot)$ consistent with the
kinetic equations.  A forward simulation starting from $\mathbf{x}_0$ selects
a single $\omega\in\Omega$ (or, stochastically, a distribution over $\Omega$).
The information content of this selection equals that of the actual cellular
trajectory, because both are governed by the same probability measure on
$\Omega$.  The simulation produces no surplus information.

\subsection{Epistemic Blindness}
\label{sec:intro-blind}

The simulation problem leads immediately to a profound biological constraint.
Suppose a cell is in a diseased state $\mathbf{x}^*_D$ differing from a
healthy state $\mathbf{x}^*_H$.  For the cell to detect its own disease it
would need an internal mechanism mapping $\mathbf{x}^*_D$ to ``diseased'' and
$\mathbf{x}^*_H$ to ``healthy''.  But the only information processing the cell
performs is forward reaction dynamics: the evaluation of $\mathbf{v}(\mathbf{x})$.
This produces a flux vector driving the next state transition, not a Boolean
classification of the current state against a stored reference.

\begin{definition}[Epistemic Blindness]
\label{def:blind}
A dynamical system $(X,\mathbf{v})$ is \emph{epistemically blind} with respect
to a partition $\{A,B\}$ of $X$ if the flow $\mathbf{v}(x)$ does not encode
membership in $A$ or $B$: there exists no function $f:TX\to\{A,B\}$ such that
$f(\mathbf{v}(x))$ correctly classifies $x$ for all $x\in X$.
\end{definition}

\begin{proposition}[Cells are epistemically blind to disease]
\label{prop:blind}
Under generic kinetic models (mass-action, Michaelis-Menten), a cell operating
via forward reaction dynamics is epistemically blind with respect to the
partition \emph{healthy}/\emph{diseased}.
\end{proposition}

\begin{proof}
The kinetic rate law $v_{ij}([C_i],[C_j],k_{ij})$ at any reaction edge depends
only on current concentrations and rate constants.  It contains no reference to
a healthy reference state $\mathbf{x}^*_H$.  At the diseased state
$\mathbf{x}^*_D$, the flux vector $\mathbf{v}(\mathbf{x}^*_D)$ encodes only
the direction of the next state transition; it carries no information about the
distance $\|\mathbf{x}^*_D - \mathbf{x}^*_H\|$ unless that distance is
directly reflected in current concentrations.  Disease states frequently
involve regulatory rewiring, post-translational modifications, or epigenetic
changes that alter the parameters of $\mathbf{v}$ rather than concentrations
directly, making the current flux indistinguishable from that of a normal cell
elsewhere in the healthy attractor.  Hence no function of
$\mathbf{v}(\mathbf{x}^*_D)$ alone reliably distinguishes diseased from
healthy.  $\square$
\end{proof}

The clinical implication is well-known: cancer cells do not self-destruct upon
becoming malignant; infected cells do not directly alert the immune system
through metabolic fluxes; misfolded protein aggregates propagate
unchecked~\cite{Frauenfelder1991,Levinthal1969}.  Epistemic blindness is not a
failure of the cell but a fundamental property of autonomous dynamical systems.

\subsection{The Circuit Approach}
\label{sec:intro-circuit}

The resolution to epistemic blindness lies in adopting the perspective of an
external observer who knows the network topology $G$ and the healthy attractor
$\mathbf{x}^*_H$.  Such an observer can determine from partial measurements
whether a given cell state lies within the healthy basin of attraction or has
escaped to a pathological attractor.  Doing so requires solving an inverse
problem: given partial observations, infer the complete state.

This paper develops a systematic framework for that inverse problem.  The core
idea is to model the reaction network as an electrical circuit graph $G=(V,E)$
in which each node $v_i\in V$ carries a potential equal to the chemical
potential $\mu_i^{\mathrm{chem}}$, shown below to equal the
information-theoretic categorical depth $\Hcat_i=-\log_2 P(\sigma_i)$ up to
universal constants.  Each directed edge $e_{ij}\in E$ carries a current equal
to the net reaction flux $J_{ij}$.  Kirchhoff's current law corresponds to
stoichiometric mass balance at steady state.  Kirchhoff's voltage law
corresponds to thermodynamic cycle consistency (the Wegscheider
conditions~\cite{Wegscheider1901}).  Node states are represented as fuzzy
membership functions encoding epistemic uncertainty, and the Kirchhoff
equations are lifted to fuzzy arithmetic via the Zadeh extension
principle~\cite{Zadeh1965}.

\subsection{Paper Organisation}
\label{sec:intro-org}

Section~\ref{sec:info} establishes information-theoretic foundations:
Shannon entropy~\cite{Shannon1948}, connection to thermodynamics via the Jaynes
maximum entropy principle~\cite{Jaynes1957}, and the categorical distance
metric.  Section~\ref{sec:circuit} develops the biochemical circuit model,
deriving circuit variables and Kirchhoff analogs from first principles.
Section~\ref{sec:fuzzy} introduces fuzzy state representation and fuzzy circuit
equations.  Section~\ref{sec:backward} defines backward trajectories, proves
their time-invariance, and connects them to attractor structure.
Section~\ref{sec:algorithm} presents the trajectory completion algorithm with
a full convergence proof.  Section~\ref{sec:time} develops the theory of
reactions as time generators via the Gillespie framework and the partition lag.
Section~\ref{sec:propagation} analyses categorical state propagation velocity
and derives the velocity ratio.  Section~\ref{sec:applications} covers
cell-state instantiation, disease detection, and drug design.
Section~\ref{sec:discussion} situates the framework in the literature and
discusses limitations.  Section~\ref{sec:conclusion} concludes.

%%=============================================================================
\section{Information-Theoretic Foundations}
\label{sec:info}
%%=============================================================================

\subsection{Shannon Entropy and Molecular States}
\label{sec:info-shannon}

We begin with the foundational result of information theory.

\begin{definition}[Shannon Entropy~\protect\cite{Shannon1948}]
\label{def:shannon}
Let $\Sigma$ be a discrete sample space with probability measure $P$.  The
\emph{Shannon entropy} of $P$ is
\begin{equation}
  H(P) = -\sum_{\sigma\in\Sigma} P(\sigma)\log_2 P(\sigma),
  \label{eq:shannon}
\end{equation}
measured in bits, with the convention $0\log_2 0\equiv 0$ by continuity.
\end{definition}

For a molecular species $i$ in a well-mixed compartment of volume $\Omega$,
the microstate $\sigma_i$ specifies the number of molecules $n_i$, equivalently
the molar concentration $[C_i]=n_i/(N_A\Omega)$.  We coarse-grain to a discrete
state space by binning concentrations into intervals of width $\delta c$
determined by the minimum detectable concentration difference (resolution of
the biological measurement).  The probability $P(\sigma_i)$ is the stationary
probability of that concentration bin under the network dynamics.

\begin{definition}[Categorical Depth]
\label{def:catdepth}
The \emph{categorical depth} of node $i$ in state $\sigma_i$ is the
self-information (surprisal) of that state:
\begin{equation}
  \Hcat_i(\sigma_i) \;=\; -\log_2 P(\sigma_i).
  \label{eq:catdepth}
\end{equation}
\end{definition}

Categorical depth measures how many binary questions are required to specify
that node $v_i$ is in state $\sigma_i$.  High $\Hcat_i$ denotes a rare,
informative state; low $\Hcat_i$ denotes a common, low-information state.
The Shannon entropy is the expected categorical depth:
$H(P)=\mathbb{E}[\Hcat_i(\sigma_i)]$.  These properties follow directly from
the standard theory of self-information~\cite{CoverThomas2006}.

\begin{figure*}[!htbp]
    \centering
    \includegraphics[width=\textwidth]{figures/panel_exp1.png}
    \caption{Categorical Depth--Chemical Potential Equivalence (Theorem~2.2).
    (a)~Categorical depth $\mathcal{H}_{\mathrm{cat}}$ vs.\ $\log_{10}$
    concentration for all 14 glycolytic species; point area is proportional to
    steady-state concentration, confirming the inverse-logarithmic relationship
    $\mathcal{H}_{\mathrm{cat}} = -\log_2(c_i / C_{\mathrm{tot}})$.
    (b)~$\mathcal{H}_{\mathrm{cat}}$ vs.\ normalised chemical potential
    $\hat{\mu}$ with the theoretical line
    $\mathcal{H}_{\mathrm{cat}} = \hat{\mu} + \kappa$ (grey); Pearson
    $r = 1.000$, constant offset $\kappa = 1.990$, confirming exact linear
    equivalence.
    (c)~Three-dimensional projection of
    $(\log_{10} c,\; \mathcal{H}_{\mathrm{cat}},\; \hat{\mu})$ showing all 14
    species lie on a single plane, as required by Theorem~2.2.
    (d)~Ranked categorical depths across all species; BPG exhibits the highest
    depth ($\mathcal{H}_{\mathrm{cat}} = 11.96$) consistent with its role as a
    transient high-energy intermediate, while ATP shows the lowest
    ($\mathcal{H}_{\mathrm{cat}} = 0.99$) reflecting its high abundance as the
    principal energy currency.}
    \label{fig:panel_exp1}
\end{figure*}

\subsection{Connection to Thermodynamics}
\label{sec:info-thermo}

The connection between information theory and thermodynamics was established by
Boltzmann~\cite{Boltzmann1877} and later placed on a variational footing by
Jaynes~\cite{Jaynes1957}.  We derive the connection explicitly for the
biochemical setting.

\begin{theorem}[Chemical Potential as Categorical Depth]
\label{thm:mu-catdepth}
For an ideal solution at thermodynamic equilibrium, the chemical potential of
species $i$ at concentration $[C_i]$ satisfies
\begin{equation}
  \mu_i^{\mathrm{chem}} = \mu_i^\circ + RT\ln[C_i]
    = -\kB T\ln 2\cdot\Hcat_i(\sigma_i) + c_i,
  \label{eq:mu-catdepth}
\end{equation}
where $R=N_A\kB$ is the gas constant, $T$ is absolute temperature, $\mu_i^\circ$
is the standard chemical potential, and
$c_i = \mu_i^\circ + RT\ln 2\cdot\log_2 Z_i$ is a node-specific constant
absorbing the partition function $Z_i$.  Thus the circuit node potential equals
the categorical depth up to the thermal energy scale $\kB T\ln 2$.
\end{theorem}

\begin{proof}
For species $i$ in an ideal solution, the Boltzmann factor gives the
equilibrium probability of microstate $\sigma_i$
as~\cite{NichollsFerguson2002,Gibbs1875}
\[
  P(\sigma_i) = Z_i^{-1}\exp\!\Bigl(-\frac{\mu_i^{\mathrm{chem}}}{RT}\Bigr),
\]
where $Z_i$ is the molecular partition function for species $i$.  Taking
$-\log_2$ of both sides:
\[
  \Hcat_i = \frac{\mu_i^{\mathrm{chem}}}{RT\ln 2} + \log_2 Z_i.
\]
Solving for $\mu_i^{\mathrm{chem}}$:
\[
  \mu_i^{\mathrm{chem}} = RT\ln 2\cdot\Hcat_i - RT\ln 2\cdot\log_2 Z_i.
\]
Writing this as $-\kB T\ln 2\cdot\Hcat_i + c_i$ with
$c_i = \mu_i^\circ + RT\ln 2\cdot\log_2 Z_i$ (where the sign is consistent
with $\mu_i^{\mathrm{chem}} = \mu_i^\circ + RT\ln[C_i]$ from statistical
mechanics: the two expressions are equal by the definition of the partition
function).  $\square$
\end{proof}

Theorem~\ref{thm:mu-catdepth} is the keystone of the circuit model: the
electrical potential at a circuit node equals, up to a universal factor
$\kB T\ln 2$, the categorical depth of the corresponding molecular species.
Reactions run downhill in categorical depth space, exactly as electric current
flows from high to low potential.

\begin{theorem}[Jaynes Maximum Entropy Characterises Equilibrium Networks]
\label{thm:jaynes}
The stationary distribution of a biochemical reaction network at thermodynamic
equilibrium is the maximum-entropy distribution over microstates subject to
the linear conservation constraints imposed by the null space of
$\mathbf{S}^T$~\cite{Jaynes1957}.
\end{theorem}

\begin{proof}
The Jaynes maximum entropy principle~\cite{Jaynes1957,CoverThomas2006} states
that the least-biased distribution satisfying constraints
$\mathbb{E}[f_k]=c_k$ ($k=1,\ldots,m$) is the Gibbs distribution
$P^*(\sigma)\propto\exp(-\sum_k\lambda_k f_k(\sigma))$.  For a reaction
network the natural constraints are the conserved moieties
$\ell_k(\mathbf{x})=\mathbf{L}_k\cdot\mathbf{x}=\mathrm{const}$, where
$\mathbf{L}_k$ spans the left null space of $\mathbf{S}$ (rows of $\mathbf{S}^T$
that are zero).  Identifying the Lagrange multipliers $\lambda_k$ with the
chemical potentials of the conserved quantities yields
$P^*(\sigma)\propto e^{-G(\sigma)/(RT)}$, precisely the Boltzmann-Gibbs
distribution~\cite{NichollsFerguson2002}.  $\square$
\end{proof}

\begin{remark}
Theorem~\ref{thm:jaynes} establishes that $\Hcat_i$ is not a formal analogy
but has genuine thermodynamic content: $\Hcat_i = G_i/(RT\ln 2) + \mathrm{const}$,
so the information content of a molecular state is proportional to its Gibbs
free energy above a reference level.  The maximum-entropy principle
simultaneously determines both the thermodynamic distribution and the
information content.
\end{remark}

\subsection{Categorical Distance}
\label{sec:info-catdist}

\begin{definition}[Categorical Distance]
\label{def:catdist}
The \emph{categorical distance} between two states $\sigma$ and $\sigma'$ of
node $v_i$ is
\begin{equation}
  \dcat(\sigma,\sigma') = \bigl|\Hcat_i(\sigma)-\Hcat_i(\sigma')\bigr|
                        = \Bigl|\log_2\frac{P(\sigma')}{P(\sigma)}\Bigr|.
  \label{eq:catdist}
\end{equation}
\end{definition}

\begin{proposition}[$\dcat$ is a metric on states of strictly positive $P$]
\label{prop:catdist-metric}
The function $\dcat:\Sigma\times\Sigma\to\mathbb{R}_{\geq 0}$ satisfies all
four metric axioms when $P(\sigma)>0$ for all distinct $\sigma\in\Sigma$.
\end{proposition}

\begin{proof}
\emph{Non-negativity:} $|\cdot|\geq 0$, so $\dcat\geq 0$.
\emph{Identity of indiscernibles:} $\dcat(\sigma,\sigma')=0
\Leftrightarrow P(\sigma)=P(\sigma')$; since $P>0$ everywhere and we take $P$
injective on distinct states, this gives $\sigma=\sigma'$.
\emph{Symmetry:} $|\Hcat(\sigma)-\Hcat(\sigma')|=|\Hcat(\sigma')-\Hcat(\sigma)|$.
\emph{Triangle inequality:} For any $\sigma,\sigma',\sigma''\in\Sigma$,
\[
  |\Hcat(\sigma)-\Hcat(\sigma'')| \leq |\Hcat(\sigma)-\Hcat(\sigma')|
  +|\Hcat(\sigma')-\Hcat(\sigma'')|
\]
by the standard absolute-value triangle inequality.  $\square$
\end{proof}

\begin{proposition}[Categorical distance and reaction free energy]
\label{prop:catdist-dG}
For a reaction converting species $i$ to species $j$, the reaction free energy
satisfies
\begin{equation}
  \Delta G_{ij} = \mu_j^{\mathrm{chem}}-\mu_i^{\mathrm{chem}}
    = -\kB T\ln 2\cdot\bigl(\Hcat_j-\Hcat_i\bigr) + (c_j-c_i).
  \label{eq:catdist-dG}
\end{equation}
The reaction is spontaneous ($\Delta G_{ij}<0$) if and only if the corrected
depth difference $\Hcat_j - \Hcat_i > (c_j-c_i)/(\kB T\ln 2)$.
\end{proposition}

\begin{proof}
Subtract the two instances of Theorem~\ref{thm:mu-catdepth}:
$\Delta G_{ij} = -\kB T\ln 2\,(\Hcat_j-\Hcat_i)+(c_j-c_i)$.
Spontaneity $\Delta G_{ij}<0$ requires
$\kB T\ln 2(\Hcat_j-\Hcat_i) > c_j-c_i$, giving the stated condition.
$\square$
\end{proof}

%%=============================================================================
\section{The Biochemical Circuit Model}
\label{sec:circuit}
%%=============================================================================

\subsection{Graph Representation}
\label{sec:circuit-graph}

\begin{definition}[Biochemical Reaction Network Graph]
\label{def:graph}
A \emph{biochemical reaction network} is a directed weighted graph
$G=(V,E,w)$ where:
\begin{itemize}[noitemsep]
  \item $V=\{v_1,\ldots,v_N\}$: species nodes---metabolites, proteins, protein
        complexes, RNA molecules, and signalling intermediates.
  \item $E\subseteq V\times V$: reaction edges.  The directed edge
        $e_{ij}=(v_i,v_j)\in E$ represents an elementary reaction step in
        which species $i$ is consumed (or acts as a reactant) and species $j$
        is produced.  Multi-substrate reactions are decomposed into binary
        edges via intermediate complex nodes following standard
        practice~\cite{HeinrichSchuster1996,Klipp2005}.
  \item $w:E\to\mathbb{R}_{>0}$: edge weight $w(e_{ij})=k_{ij}$, the
        elementary rate constant.
\end{itemize}
Each node $v_i$ carries a state variable $[C_i](t)\in\mathbb{R}_{\geq 0}$
representing the molar concentration of species $i$ at time $t$.
\end{definition}

\begin{definition}[Stoichiometric Matrix]
\label{def:stoich}
The \emph{stoichiometric matrix} $\mathbf{S}\in\mathbb{Z}^{N\times R}$ has
entries $S_{ir}$ equal to the net stoichiometric coefficient of species $i$ in
reaction $r$ ($S_{ir}>0$: produced; $S_{ir}<0$: consumed; $S_{ir}=0$:
absent).  The system ODE follows as $\dot{\mathbf{x}}=\mathbf{S}\mathbf{v}(\mathbf{x})$
(equation~\ref{eq:ode}).  In the circuit analogy, $\mathbf{S}$ plays the role
of the node-edge incidence matrix.
\end{definition}

The graph $G$ encodes all possible molecular interactions in the cell.  Its
topology determines which states are reachable, which conservation laws hold
(null space of $\mathbf{S}^T$), and which thermodynamic cycles exist (cycle
space of $G$).

\subsection{Circuit Variables}
\label{sec:circuit-vars}

\begin{definition}[Node Potential]
\label{def:potential}
The \emph{circuit potential} at node $v_i$ is the chemical potential:
\begin{equation}
  \phi_i \;=\; \mu_i^{\mathrm{chem}} \;=\; \mu_i^\circ + RT\ln[C_i].
  \label{eq:potential}
\end{equation}
By Theorem~\ref{thm:mu-catdepth},
$\phi_i = -\kB T\ln 2\cdot\Hcat_i + c_i$.
\end{definition}

\begin{definition}[Edge Current]
\label{def:current}
The \emph{circuit current} on edge $e_{ij}$ is the net reaction flux:
\begin{equation}
  J_{ij} = v_{ij}^+ - v_{ij}^- = k_{ij}[C_i] - k_{ji}[C_j],
  \label{eq:current}
\end{equation}
where $v_{ij}^+=k_{ij}[C_i]$ is the forward flux and
$v_{ij}^-=k_{ji}[C_j]$ is the reverse flux.  The existence of a nonzero
reverse rate $k_{ji}$ is required by microscopic
reversibility~\cite{Wegscheider1901,Lewis1925,Onsager1931}.
\end{definition}

\begin{definition}[Edge Conductance]
\label{def:conductance}
The \emph{circuit conductance} of edge $e_{ij}$ is
\begin{equation}
  \Gcond_{ij} = \frac{k_{ij}[C_i]}{RT},
  \label{eq:conductance}
\end{equation}
in units of mol\,L$^{-1}$\,J$^{-1}$\,s$^{-1}$.
\end{definition}

Table~\ref{tab:analogy} summarises the complete circuit--biochemistry
correspondence.

\begin{table}[t]
\caption{Complete circuit--biochemistry correspondence.  $k_{ij}$: forward
rate constant; $[C_i]$: concentration; $R$: gas constant; $T$: temperature.}
\label{tab:analogy}
\small
\begin{tabular}{@{}lp{3.9cm}@{}}
\toprule
\textbf{Circuit element} & \textbf{Biochemical equivalent} \\
\midrule
Node potential $\phi_i$      & Chemical potential $\mu_i^{\mathrm{chem}}$ \\
Potential difference         & Reaction $\Delta G_{ij}$ \\
Edge current $J_{ij}$        & Net reaction flux \\
Conductance $\Gcond_{ij}$    & $k_{ij}[C_i]/(RT)$ \\
Resistance $\mathcal{R}_{ij}$& $RT/(k_{ij}[C_i])$ \\
EMF source                   & Standard free energy $\Delta G^\circ_{ij}$ \\
Capacitance at $v_i$         & Buffer capacity $d[C_i]/d\phi_i$ \\
Ground ($\phi=0$)             & Reference state ($\mu^\circ$ basis) \\
Ohm's law $J=\Gcond\Delta\phi$ & Near-equilibrium mass action \\
Diode / transistor           & Saturable enzyme \\
\bottomrule
\end{tabular}
\end{table}

\subsection{Kirchhoff Analogs}
\label{sec:circuit-kirchhoff}

\begin{theorem}[Biochemical KCL: Stoichiometric Mass Balance]
\label{thm:kcl}
At steady state, the net flux entering any node $v_i$ vanishes:
\begin{equation}
  \sum_{\substack{j:\\(j,i)\in E}} J_{ji}
  - \sum_{\substack{j:\\(i,j)\in E}} J_{ij} = 0,
  \qquad \forall\,v_i\in V.
  \label{eq:kcl}
\end{equation}
This is the exact biochemical analog of Kirchhoff's current
law~\cite{Kirchhoff1845}.
\end{theorem}

\begin{proof}
At steady state, $d[C_i]/dt=0$.  Expanding:
$\sum_r S_{ir}v_r(\mathbf{x})=0$.
Decomposing each net flux $v_r=v_r^+-v_r^-$ and associating each directed
elementary step with an edge $e_{ij}$ gives
$\sum_{(j,i)\in E}J_{ji}=\sum_{(i,j)\in E}J_{ij}$,
which is~\eqref{eq:kcl}.  The structure is identical to KCL:
sum of currents in equals sum of currents out at every node~\cite{Kirchhoff1845}.
$\square$
\end{proof}

\begin{theorem}[Biochemical KVL: Thermodynamic Cycle Consistency]
\label{thm:kvl}
Around any directed cycle
$\mathcal{C}=(v_{i_1},v_{i_2},\ldots,v_{i_k},v_{i_1})$ in $G$, the sum of
potential differences vanishes:
\begin{equation}
  \sum_{(i,j)\in\mathcal{C}} \Delta\phi_{ij} = 0,
  \label{eq:kvl}
\end{equation}
where $\Delta\phi_{ij}=\phi_j-\phi_i=RT\ln([C_j]/[C_i])+\Delta\mu^\circ_{ij}$.
This is the biochemical analog of Kirchhoff's voltage law~\cite{Kirchhoff1845}
and is equivalent to the Wegscheider conditions~\cite{Wegscheider1901}.
\end{theorem}

\begin{proof}
The sum telescopes:
\[
\sum_{(i,j)\in\mathcal{C}}(\phi_j-\phi_i)
=\phi_{i_2}-\phi_{i_1}+\phi_{i_3}-\phi_{i_2}+\cdots+\phi_{i_1}-\phi_{i_k}=0.
\]
The concentration-dependent part gives
$RT\sum_{(i,j)}\ln([C_j]/[C_i])=RT\ln\prod_{(i,j)}[C_j]/[C_i]=RT\ln 1=0$
(the product over a closed cycle telescopes to 1).  The standard-state part
$\sum_{(i,j)}\Delta\mu^\circ_{ij}=0$ is the Wegscheider
condition~\cite{Wegscheider1901,Schuster1991}, which is required for
thermodynamic consistency.  At equilibrium, microscopic reversibility further
gives $\prod_{(i,j)\in\mathcal{C}}k_{ij}/k_{ji}=1$, the kinetic form of the
Wegscheider condition.  The structure is identical to KVL:
$\sum_{\mathrm{loop}}\Delta V=0$~\cite{Kirchhoff1845}.  $\square$
\end{proof}

\begin{corollary}[Ohm's Law Analog]
\label{cor:ohm}
For mass-action kinetics near thermodynamic equilibrium, the net flux satisfies
\begin{equation}
  J_{ij}\;\approx\;\Gcond_{ij}\,\Delta\phi_{ij},
\end{equation}
in exact formal correspondence with Ohm's law $I=G\cdot V$.
\end{corollary}

\begin{proof}
Net flux $J_{ij}=k_{ij}[C_i]-k_{ji}[C_j]$.
From detailed balance $k_{ji}=k_{ij}e^{-\Delta\mu^\circ_{ij}/(RT)}$.
Writing $\Delta\phi_{ij}=RT\ln([C_j]/[C_i])+\Delta\mu^\circ_{ij}$,
one obtains $J_{ij}=k_{ij}[C_i](1-e^{-\Delta\phi_{ij}/(RT)})$.
For $|\Delta\phi_{ij}|\ll RT$, linearising the exponential:
$J_{ij}\approx(k_{ij}[C_i]/(RT))\Delta\phi_{ij}=\Gcond_{ij}\Delta\phi_{ij}$.
This is the near-equilibrium form of the linear Onsager relations~\cite{Onsager1931}.
$\square$
\end{proof}

\subsection{Edge Conductance from Kinetics}
\label{sec:circuit-conductance}

\begin{definition}[Eyring--Arrhenius Conductance]
\label{def:arrhenius-cond}
For an elementary reaction step with Gibbs activation free energy
$\Delta G^\ddagger_{ij}$, the Eyring transition-state
rate~\cite{Eyring1935,Evans1935} is
\begin{equation}
  k_{ij} = \frac{\kB T}{h}\exp\!\left(-\frac{\Delta G^\ddagger_{ij}}{RT}\right),
  \label{eq:eyring}
\end{equation}
where $h$ is Planck's constant.  The Arrhenius approximation~\cite{Arrhenius1889}
identifies $\Delta G^\ddagger\approx E_a$ (activation energy) and
$A_{ij}=\kB T/h$ (pre-exponential factor).  Substituting
into~\eqref{eq:conductance}:
\begin{equation}
  \Gcond_{ij}^{\mathrm{Arr}} = \frac{[C_i]\kB T}{hRT}
    \exp\!\left(-\frac{\Delta G^\ddagger_{ij}}{RT}\right).
  \label{eq:cond-arr}
\end{equation}
Conductance decreases exponentially with the activation barrier.
\end{definition}

\begin{definition}[Michaelis-Menten Conductance]
\label{def:mm-cond}
For an enzymatic reaction $S+E\rightleftharpoons ES\to P+E$ at quasi-steady
state~\cite{MichaelisMenten1913,Briggs1925}, with $v_{\max}=k_{\mathrm{cat}}[E]_T$
and Michaelis constant $K_m=(k_{-1}+k_{\mathrm{cat}})/k_1$, the effective
conductance between substrate node $S$ and product node $P$ is
\begin{equation}
  \Gcond_{SP}^{\mathrm{MM}} = \frac{k_{\mathrm{cat}}[E]_T}{RT(K_m+[S])}.
  \label{eq:cond-mm}
\end{equation}
At saturation $([S]\gg K_m)$:
$\Gcond_{SP}^{\mathrm{MM}}\to k_{\mathrm{cat}}[E]_T/(RT)$ (current-source regime).
At subsaturation $([S]\ll K_m)$:
$\Gcond_{SP}^{\mathrm{MM}}\approx(k_{\mathrm{cat}}/K_m)[E]_T[S]/(RT)$
(ohmic regime, proportional to $[S]$).
\end{definition}

\begin{remark}
The Michaelis-Menten enzyme functions as a nonlinear circuit element: ohmic at
low substrate (behaving like a resistor) and saturating at high substrate
(behaving like a current source).  This is the biochemical transistor.  The
transition at $[S]\approx K_m$ introduces the nonlinearity responsible for
bistability, sustained oscillations, and ultrasensitivity in cellular signalling
networks~\cite{Tyson2003,HeinrichSchuster1996,Strogatz1994}.
\end{remark}

\begin{figure*}[!htbp]
    \centering
    \includegraphics[width=\textwidth]{figures/panel_exp2.png}
    \caption{Kirchhoff Current Law Validation at Metabolic Steady State
    (Theorem~3.3).
    (a)~Steady-state concentrations for the 10-node glycolytic circuit after
    numerical integration to convergence; glucose (GLC) dominates at
    ${\sim}1.0\;\mathrm{mM}$ while 1,3-BPG is the rarest species at
    ${\sim}1\;\mu\mathrm{M}$.
    (b)~Absolute net categorical flux $|J_{\mathrm{net}}|$ at each node on a
    logarithmic scale; all residuals fall below $1.7 \times 10^{-6}\;\mathrm{mM\,s^{-1}}$,
    confirming that $\sum_j J_{ij} \approx 0$ at every node (Kirchhoff's
    current law in metabolic form).
    (c)~Three-dimensional visualisation of enzyme flux magnitude and
    steady-state concentration across the nine reaction steps, illustrating
    uniform flux (${\sim}0.100\;\mathrm{mM\,s^{-1}}$) through a heterogeneous
    concentration landscape.
    (d)~Enzyme-resolved flux profile showing all nine glycolytic enzymes
    (HK through PK) carry effectively identical categorical current at steady
    state, with deviations $< 0.002\%$.}
    \label{fig:panel_exp2}
\end{figure*}

%%=============================================================================
\section{Fuzzy State Representation}
\label{sec:fuzzy}
%%=============================================================================

\subsection{Zadeh Fuzzy Sets Applied to Node States}
\label{sec:fuzzy-zadeh}

In practice the state of any node $v_i$ is never perfectly known.  Omics
measurements have finite precision; many species are unobserved; and intrinsic
stochasticity broadens the distribution over copy
numbers~\cite{Elowitz2002,vanKampen2007}.  We represent this epistemic
uncertainty using fuzzy sets~\cite{Zadeh1965}.

\begin{definition}[Fuzzy Node State]
\label{def:fuzzy-state}
A \emph{fuzzy state} of node $v_i$ is a membership function
$\tilde{\mu}_i:\mathbb{R}_{\geq 0}\to[0,1]$, where $\tilde{\mu}_i(c)$ is the
\emph{degree of membership} of concentration $c$ in the fuzzy set representing
the current state.  $\tilde{\mu}_i(c)=1$ means $c$ is fully consistent with
all available information; $\tilde{\mu}_i(c)=0$ means $c$ is excluded.
\end{definition}

\begin{definition}[Support, Core, and $\alpha$-cuts]
\label{def:alpha-cut}
For fuzzy state $\tilde{\mu}_i$:
\begin{align}
  \mathrm{supp}(\tilde{\mu}_i) &= \{c\geq 0:\tilde{\mu}_i(c)>0\}, \\
  \mathrm{core}(\tilde{\mu}_i) &= \{c\geq 0:\tilde{\mu}_i(c)=1\}, \\
  [\tilde{\mu}_i]^\alpha        &= \{c\geq 0:\tilde{\mu}_i(c)\geq\alpha\},
  \quad\alpha\in(0,1].
\end{align}
When $\tilde{\mu}_i$ is upper semi-continuous, $[\tilde{\mu}_i]^\alpha$ is a
closed interval $[\underline{c}_i^\alpha,\overline{c}_i^\alpha]$
\cite{Zadeh1978,Zimmermann2001book}.
\end{definition}

\begin{definition}[Trapezoidal and Triangular Membership Functions]
\label{def:tri-trap}
A \emph{trapezoidal} fuzzy state with parameters $a\leq b\leq d\leq e$:
\begin{equation}
  \tilde{\mu}_i^{\mathrm{trap}}(c) =
  \begin{cases}
    (c-a)/(b-a) & a\leq c < b \\
    1           & b\leq c\leq d \\
    (e-c)/(e-d) & d< c\leq e \\
    0           & \text{otherwise.}
  \end{cases}
  \label{eq:trapezoid}
\end{equation}
A \emph{triangular} membership function is the special case $b=d$.  Trapezoidal
functions arise from interval measurements with error bounds; triangular
functions from point estimates with uncertainty.
\end{definition}

\begin{definition}[Fuzzy Arithmetic via the Zadeh Extension Principle]
\label{def:extension}
The Zadeh extension principle~\cite{Zadeh1965,Zadeh1978} lifts a crisp function
$f:\mathbb{R}^n\to\mathbb{R}$ to fuzzy operands
$\tilde{\mu}_1,\ldots,\tilde{\mu}_n$:
\begin{equation}
  \tilde{\mu}_f(z) = \sup_{\substack{c_1,\ldots,c_n\geq 0\\f(\mathbf{c})=z}}
    \min_{k}\tilde{\mu}_k(c_k).
  \label{eq:extension}
\end{equation}
For $\alpha$-cuts (closed intervals), this reduces to interval
arithmetic~\cite{Zimmermann2001book}:
\begin{align}
  [\underline{a},\overline{a}]+[\underline{b},\overline{b}]
    &=[\underline{a}+\underline{b},\,\overline{a}+\overline{b}], \label{eq:ia-add}\\
  [\underline{a},\overline{a}]-[\underline{b},\overline{b}]
    &=[\underline{a}-\overline{b},\,\overline{a}-\underline{b}], \label{eq:ia-sub}\\
  [\underline{a},\overline{a}]\cdot[\underline{b},\overline{b}]
    &=\bigl[\min_{p,q}\{pq\},\max_{p,q}\{pq\}\bigr],\quad p\in[\underline{a},\overline{a}],\;
      q\in[\underline{b},\overline{b}], \label{eq:ia-mul}\\
  [\underline{a},\overline{a}]/[\underline{b},\overline{b}]
    &=[\underline{a},\overline{a}]\cdot[1/\overline{b},\,1/\underline{b}],
    \quad 0\notin[\underline{b},\overline{b}]. \label{eq:ia-div}
\end{align}
\end{definition}

\subsection{Physical Interpretation of Fuzziness}
\label{sec:fuzzy-physical}

Fuzzy node states encode two distinct but complementary sources of uncertainty.

\emph{(a) Measurement uncertainty.}  Analytical techniques (mass spectrometry,
fluorescence microscopy, RNA-seq) yield a measured value with a precision
interval $[C_i-\epsilon,C_i+\epsilon]$.  This maps directly to a trapezoidal
membership function with flat top equal to the measurement interval.

\emph{(b) Intrinsic stochastic fluctuations.}  In a mammalian cell
(volume $\sim 10^{-12}$\,L), a nanomolar concentration corresponds to
$\sim 600$ molecules; at picomolar levels, this drops to $<1$ molecule on
average.  The probability distribution over copy numbers in the stationary
Chemical Master Equation~\cite{vanKampen2007,Gillespie1977} is then genuinely
broad.  The membership function $\tilde{\mu}_i(c)$ can be set proportional to
the stationary CME probability density, giving a probabilistic fuzzy state.

\begin{remark}
Probability theory and fuzzy set theory are complementary.  A probability
distribution answers: \emph{which} concentration value is the true one?  A
fuzzy membership function answers: \emph{how well} does a given concentration
satisfy the node's biochemical role?  Both are needed in trajectory completion:
probability distributions govern inference; membership functions encode
role-based filtering (e.g., ``enzyme is in the active concentration range'').
\end{remark}

\subsection{Fuzzy Circuit Equations}
\label{sec:fuzzy-circuit}

\begin{definition}[Fuzzy Network State and Hausdorff Metric]
\label{def:fuzzy-net}
The \emph{fuzzy network state} is the tuple
$\widetilde{\mathbf{X}}=(\tilde{\mu}_1,\ldots,\tilde{\mu}_N)
\in\widetilde{\mathcal{X}}=\prod_{i=1}^N\mathcal{F}(\mathbb{R}_{\geq 0})$,
where $\mathcal{F}(\mathbb{R}_{\geq 0})$ denotes the set of upper
semi-continuous normal fuzzy sets on $\mathbb{R}_{\geq 0}$.  The
\emph{Hausdorff--Pompeiu metric} on $\mathcal{F}(\mathbb{R}_{\geq 0})$
is~\cite{Hausdorff1914}
\begin{equation}
  d_H(\tilde{\mu},\tilde{\nu})=\sup_{\alpha\in(0,1]}
    d_{\mathrm{H}}\bigl([\tilde{\mu}]^\alpha,[\tilde{\nu}]^\alpha\bigr),
  \label{eq:hausdorff}
\end{equation}
where for closed intervals
$d_{\mathrm{H}}([a,b],[c,d])=\max(|a-c|,|b-d|)$.  The product metric on
$\widetilde{\mathcal{X}}$ is
$d(\widetilde{\mathbf{X}},\widetilde{\mathbf{Y}})=\max_i d_H(\tilde{\mu}_i,\tilde{\nu}_i)$.
\end{definition}

\begin{theorem}[Fuzzy KCL]
\label{thm:fuzzy-kcl}
The Kirchhoff current law~\eqref{eq:kcl} lifts to fuzzy states via the
extension principle.  At each node $v_i$, the fuzzy flux balance constraint
$\widetilde{J}_i^{\mathrm{in}}=\widetilde{J}_i^{\mathrm{out}}$ is enforced by
intersecting the current membership function $\tilde{\mu}_i$ with the set of
concentrations consistent with the fuzzy mass balance:
\begin{equation}
  \tilde{\mu}_i \;\leftarrow\; \tilde{\mu}_i \cap \tilde{\mu}_i^{\mathrm{KCL}},
  \label{eq:fuzzy-kcl}
\end{equation}
where $(\tilde{\mu}_1\cap\tilde{\mu}_2)(c)=\min(\tilde{\mu}_1(c),\tilde{\mu}_2(c))$
is the standard fuzzy intersection~\cite{Zadeh1965}.
\end{theorem}

\begin{proof}
Apply Definition~\ref{def:extension} to the crisp constraint
$\sum_jk_{ji}[C_j]=\sum_jk_{ij}[C_i]$ (Theorem~\ref{thm:kcl}) viewed as a
function of the concentrations.  At each $\alpha$-level, the fuzzy intervals
for all concentrations are propagated through the linear flux equations by
interval arithmetic~\eqref{eq:ia-add}--\eqref{eq:ia-mul}.  The resulting
interval for $[C_i]$ consistent with fuzzy KCL at level $\alpha$ is
$[\underline{c}_i^{\mathrm{KCL},\alpha},\overline{c}_i^{\mathrm{KCL},\alpha}]$,
defining $[\tilde{\mu}_i^{\mathrm{KCL}}]^\alpha$.  Assembling over all $\alpha$
gives $\tilde{\mu}_i^{\mathrm{KCL}}$, and the intersection
$\tilde{\mu}_i\cap\tilde{\mu}_i^{\mathrm{KCL}}$ restricts the state of $v_i$
to concentrations simultaneously consistent with measurement and flux balance.
$\square$
\end{proof}

\begin{theorem}[Fuzzy KVL]
\label{thm:fuzzy-kvl}
For each directed cycle $\mathcal{C}$ in $G$, the thermodynamic cycle
constraint~\eqref{eq:kvl} in fuzzy form requires:
\begin{equation}
  0\;\in\;\mathrm{core}\!\left(\widetilde{\Phi}_{\mathcal{C}}\right),
  \label{eq:fuzzy-kvl}
\end{equation}
where $\widetilde{\Phi}_{\mathcal{C}}$ is the fuzzy real number obtained by
applying the extension principle to
$\sum_{(i,j)\in\mathcal{C}}(\phi_j-\phi_i)$.  Any assignment violating this
condition has its membership reduced.
\end{theorem}

\begin{proof}
By Theorem~\ref{thm:kvl}, in the crisp case the cycle potential sum is zero
exactly.  Applying the extension principle to the sum
$\sum_{(i,j)\in\mathcal{C}}(RT\ln[C_j]-RT\ln[C_i]+\Delta\mu^\circ_{ij})$
produces the fuzzy number $\widetilde{\Phi}_{\mathcal{C}}$.  The constraint
$\widetilde{\Phi}_{\mathcal{C}}(0)=1$ (zero must be in the core) is enforced
by back-propagating through the interval arithmetic: at each $\alpha$-level,
the intervals of cycle-node concentrations are restricted so that
$[\widetilde{\Phi}_{\mathcal{C}}]^\alpha$ contains zero.  $\square$
\end{proof}

\begin{figure*}[!htbp]
    \centering
    \includegraphics[width=\textwidth]{figures/panel_exp4.png}
    \caption{Fuzzy State Completion via Kirchhoff Constraints (Theorem~6.2).
    (a)~Grouped bars comparing true steady-state concentrations (coloured) with
    na\"{\i}ve uniform prior ($0.1\;\mathrm{mM}$, grey) for the eight
    unobserved glycolytic intermediates; the prior deviates by up to
    $99\times$ (BPG).
    (b)~Completed concentrations vs.\ true values for all eight species;
    points lie on the identity line ($y = x$) with mean absolute relative
    error $\mathrm{MARE} = 2.1 \times 10^{-10}$, demonstrating that the
    Kirchhoff completion operator recovers the full state from partial
    observations with effectively zero error.
    (c)~Three-dimensional view of (species index, true concentration,
    completed concentration) confirming pointwise agreement across the full
    metabolite vector.
    (d)~Log-scale comparison of relative reconstruction error: na\"{\i}ve
    prior (orange, $\mathrm{MARE} = 14.7$) vs.\ Kirchhoff completion (blue,
    $\mathrm{MARE} \sim 10^{-10}$); Monte Carlo robustness testing over 100
    trials with $5\%$ observation noise yields $100\%$ of completions within
    $20\%$ of the true value ($\mathrm{MC\text{-}MARE} = 0.0046$).}
    \label{fig:panel_exp4}
\end{figure*}

%%=============================================================================
\section{Backward Trajectories and Temporal Invariants}
\label{sec:backward}
%%=============================================================================

\subsection{Formal Definition}
\label{sec:backward-def}

\begin{definition}[Reaction Network as a CTMC]
\label{def:markov}
For reaction network $G=(V,E,w)$ with mass-action kinetics, the stochastic
dynamics of the molecule-count vector $\mathbf{n}(t)\in\mathbb{Z}_{\geq 0}^N$
form a continuous-time Markov chain (CTMC) with propensity functions
\cite{Gillespie1977,vanKampen2007}:
\begin{equation}
  a_{ij}(\mathbf{n}) = k_{ij}\,\frac{n_i!}{(n_i-\nu^-_{ij})!}\,
    \Omega^{1-\nu^-_{ij}},
  \label{eq:propensity}
\end{equation}
where $\nu^-_{ij}$ is the molecularity of species $i$ as reactant in $e_{ij}$,
$\Omega$ is the compartment volume in units of the reference volume, and
$a_{ij}(\mathbf{n})$ is the instantaneous probability per unit time of
reaction $e_{ij}$ firing.
\end{definition}

\begin{definition}[MAP Backward Trajectory]
\label{def:backward-traj}
Let node $v_i$ be observed in state $\sigma_i$ at time $T$.  The
\emph{backward trajectory} $\Back_i$ is the maximum-a-posteriori (MAP) causal
history of that observation:
\begin{equation}
  \Back_i = \argmax_{\substack{\mathbf{x}(\cdot)\,:\,x_i(T)=\sigma_i\\
    \mathbf{x}(\cdot)\,\text{consistent with}\,G}}
  P\bigl(\mathbf{x}(\cdot)\mid x_i(T)=\sigma_i,\,G\bigr).
  \label{eq:backward-traj}
\end{equation}
It is the most probable trajectory, consistent with the network kinetics and
topology, that terminates in state $\sigma_i$ at time $T$.
\end{definition}

Computationally, $\Back_i$ is obtained by the Viterbi
algorithm~\cite{Viterbi1967,Forney1973} applied to the CTMC: starting from
$\sigma_i$ at time $T$, trace backward through the reaction graph selecting
at each step the predecessor reaction with the highest conditional probability.

\begin{definition}[Categorical Address]
\label{def:cat-address}
The \emph{categorical address} of node $v_i$ in state $\sigma_i$ is the
backward trajectory $\Back_i(\sigma_i,G)$.  It is the time-invariant identity
of that state in the network's state space.
\end{definition}

\subsection{Time-Invariance Theorem}
\label{sec:backward-timeinv}

\begin{theorem}[Time-Invariance of Backward Trajectories]
\label{thm:timeinv}
For an autonomous (time-homogeneous) reaction network $G$ with time-independent
rate constants, the MAP backward trajectory $\Back_i$ depends only on the
current state $\sigma_i$ and the network $G$, not on the absolute observation
time $T$:
\begin{equation}
  \Back_i = \Back_i(\sigma_i,\,G),\quad\text{independent of }T.
  \label{eq:timeinv}
\end{equation}
\end{theorem}

\begin{proof}
The propensity functions $a_{ij}(\mathbf{n})$ in~\eqref{eq:propensity} depend
only on the current state $\mathbf{n}$, not on $t$.  The CTMC is therefore
time-homogeneous: for all $t,s\geq 0$,
$P(\mathbf{n}(t+s)=\mathbf{m}\mid\mathbf{n}(t)=\mathbf{n})=
P(\mathbf{n}(s)=\mathbf{m}\mid\mathbf{n}(0)=\mathbf{n})$~\cite{vanKampen2007}.

Applying this shift to the conditional distribution of the past trajectory:
for any $\tau>0$,
\[
  P\bigl(\mathbf{x}(T-\tau)=\mathbf{m}\mid x_i(T)=\sigma_i,\,G\bigr)
  = P\bigl(\mathbf{x}(-\tau)=\mathbf{m}\mid x_i(0)=\sigma_i,\,G\bigr),
\]
since the CTMC statistics are identical after the time shift $T$.  The MAP
over the left-hand side is therefore identical to the MAP over the right-hand
side for all $T$, establishing~\eqref{eq:timeinv}.  $\square$
\end{proof}

\begin{corollary}[Categorical address as geometric invariant]
\label{cor:geom-inv}
The categorical address $\Back_i(\sigma_i,G)$ is a fixed geometric property
of state $\sigma_i$ within the network's state space, independent of when the
cell is observed.
\end{corollary}

\begin{proof}
Immediate from Theorem~\ref{thm:timeinv}: since $\Back_i$ is independent of
$T$, it depends only on $(\sigma_i,G)$, which are static data.  $\square$
\end{proof}

\subsection{The Categorical Address}
\label{sec:backward-address}

\begin{proposition}[Categorical address encodes network topology]
\label{prop:cat-addr-topo}
Two nodes $v_i$ and $v_j$ with the same observed concentration
($[C_i]=[C_j]$) but different topological positions in $G$ have distinct
categorical addresses: $\Back_i(\sigma,G)\neq\Back_j(\sigma,G)$ in general.
\end{proposition}

\begin{proof}
The Viterbi backward path from node $v_i$ at time $T$ selects, at each step,
the predecessor reaction with the highest propensity.  The available
predecessors are the in-edges of $v_i$ in the graph $G$.  Since $v_i$ and
$v_j$ occupy different positions in $G$, their in-edge sets are different
(by assumption of distinct topology), so the most-probable predecessor
sequence diverges, yielding distinct categorical addresses.  $\square$
\end{proof}

This proposition shows that the categorical address is a richer descriptor
than concentration alone: two molecules at the same concentration but in
different network positions are distinguished by their addresses.  This is
precisely the information that forward simulation fails to provide but that
trajectory completion recovers.

\subsection{Relationship to Attractor Structure}
\label{sec:backward-attractor}

\begin{definition}[Healthy and Diseased Attractors]
\label{def:attractors}
Let~\eqref{eq:ode} possess multiple attractors $\mathcal{A}_H$ (healthy) and
$\mathcal{A}_D$ (diseased), with basins $\mathcal{B}_H$ and $\mathcal{B}_D$.
A cell state $\mathbf{x}^*$ is \emph{healthy} if
$\mathbf{x}^*\in\mathcal{B}_H$, and \emph{diseased} if
$\mathbf{x}^*\in\mathcal{B}_D\setminus\mathcal{B}_H$.
\end{definition}

\begin{theorem}[Health status from backward trajectory convergence]
\label{thm:health-address}
A cell is healthy if and only if the backward trajectory $\Back_i(\sigma_i,G)$
of every node $v_i$ converges to a state in the healthy attractor
$\mathcal{A}_H$.  Disease corresponds to at least one node whose backward
trajectory instead converges to the diseased attractor $\mathcal{A}_D$.
\end{theorem}

\begin{proof}
\emph{Healthy $\Rightarrow$ all backward trajectories converge to $\mathcal{A}_H$}:
If $\mathbf{x}^*\in\mathcal{B}_H$, the forward trajectory converges to
$\mathcal{A}_H$ as $t\to\infty$.  The MAP backward trajectory traces the
time-reversed most probable path; for a gradient-like flow on $\mathcal{B}_H$,
this path also lies within $\mathcal{B}_H$ and converges to $\mathcal{A}_H$
as backward time $\tau\to\infty$.

\emph{Diseased $\Rightarrow$ some backward trajectory escapes to $\mathcal{A}_D$}:
If $\mathbf{x}^*\in\mathcal{B}_D$, the forward trajectory converges to
$\mathcal{A}_D$.  The MAP backward trajectory must originate from a region
consistent with this forward convergence, placing it within $\mathcal{B}_D$
and converging to $\mathcal{A}_D$.  By Proposition~\ref{prop:blind}, the cell
cannot detect this; an external observer can, by computing the categorical
address and comparing to the healthy reference.  $\square$
\end{proof}

Figure~\ref{fig:attractors} illustrates attractor basins and backward
trajectory divergence.

\begin{figure*}[!htbp]
    \centering
    \includegraphics[width=\textwidth]{figures/panel_exp3.png}
    \caption{Backward Trajectory Time-Invariance (Theorem~5.2).
    (a)~Forward-trajectory snapshots in the GLC--PYR phase plane at
    $t = 50$, $100$, and $150\;\mathrm{s}$; GLC decays monotonically while PYR
    accumulates, yielding three distinct starting states for backward
    integration.
    (b)~Cosine similarity between backward-trajectory direction vectors
    launched from the three time points; all values exceed $0.9998$,
    confirming that the direction field is time-invariant to
    $< 0.02\%$.
    (c)~Three-dimensional backward trajectories in
    (GLC, F6P, PEP) space: circles mark forward-trajectory starting states,
    stars mark backward endpoints after $20\;\mathrm{s}$ reverse integration;
    connecting lines show that all three trajectories converge to a common
    region of partition space despite originating from different time points.
    (d)~Root-mean-square deviation (RMSD) between backward endpoints as a
    function of query time; RMSD remains below $0.16\;\mathrm{mM}$ across the
    full $100\;\mathrm{s}$ window, demonstrating geometric stability of the
    backward attractor.}
    \label{fig:panel_exp3}
\end{figure*}

%%=============================================================================
\section{Trajectory Completion Algorithm}
\label{sec:algorithm}
%%=============================================================================

\subsection{Algorithm Statement}
\label{sec:alg-statement}

The trajectory completion problem takes as input:
\begin{itemize}[noitemsep]
  \item a network $G=(V,E,w)$ with standard potentials $\{\mu_i^\circ\}$;
  \item partial observations
        $\mathcal{O}=\{(v_{i_k},\sigma_{i_k})\}_{k=1}^M$, $M\leq N$;
  \item healthy reference backward trajectories $\{\Back_i^H\}$ precomputed
        from a healthy reference network state.
\end{itemize}
It outputs a fuzzy network state $\widetilde{\mathbf{X}}^*$ maximally
consistent with all constraints.

\begin{definition}[Constraint Operators]
\label{def:operators}
Define three operators on $(\widetilde{\mathcal{X}},d)$:
\begin{enumerate}[label=(\roman*),noitemsep]
  \item $\mathcal{T}_{\mathrm{KCL}}$: apply Theorem~\ref{thm:fuzzy-kcl} at
        every node, enforcing fuzzy mass balance.
  \item $\mathcal{T}_{\mathrm{KVL}}$: apply Theorem~\ref{thm:fuzzy-kvl} to
        every independent cycle in $G$ (circuit rank $R-N+1$ cycles
        suffice~\cite{HeinrichSchuster1996}).
  \item $\mathcal{T}_{\mathrm{Back}}$: for each node $v_i$, compute the
        Viterbi MAP backward trajectory~\cite{Viterbi1967,Forney1973} from
        the current most-probable concentration; restrict $\tilde{\mu}_i$ to
        concentrations whose backward trajectory is consistent with $\Back_i^H$.
\end{enumerate}
The combined operator is
$\mathcal{T}=\mathcal{T}_{\mathrm{Back}}\circ
\mathcal{T}_{\mathrm{KVL}}\circ\mathcal{T}_{\mathrm{KCL}}$.
\end{definition}

\begin{algorithm}[t]
\caption{Trajectory Completion Algorithm}
\label{alg:tc}
\begin{algorithmic}[1]
\Require Network $G$; observations $\mathcal{O}$; potentials $\{\mu_i^\circ\}$;
         reference trajectories $\{\Back_i^H\}$; tolerance $\varepsilon>0$.
\Ensure  Fuzzy state $\widetilde{\mathbf{X}}^*$.
\State \textbf{Initialise}: For observed nodes $v_{i_k}$, set
       $\tilde{\mu}_{i_k}\leftarrow\tilde{\mu}^{\mathrm{trap}}(\cdot;\,
       \sigma_{i_k}-\epsilon,\sigma_{i_k}-\epsilon/2,\sigma_{i_k}+\epsilon/2,
       \sigma_{i_k}+\epsilon)$.
       For unobserved nodes, set
       $\tilde{\mu}_i\leftarrow\mathbf{1}$ on $[c_{\min}^i,c_{\max}^i]$
       (maximum-entropy prior from database ranges).
\Repeat
  \State $\widetilde{\mathbf{X}}^{\mathrm{prev}}\leftarrow\widetilde{\mathbf{X}}$
  \For{each node $v_i\in V$}
    \State Apply $\mathcal{T}_{\mathrm{KCL}}$ at $v_i$: compute the fuzzy
           in-flux and out-flux from current neighbour states via fuzzy
           arithmetic; intersect $\tilde{\mu}_i$ with the KCL-consistent set
           (Theorem~\ref{thm:fuzzy-kcl}).
  \EndFor
  \For{each independent cycle $\mathcal{C}$ in $G$}
    \State Apply $\mathcal{T}_{\mathrm{KVL}}$: compute the fuzzy cycle
           potential; restrict node memberships to enforce
           $0\in\mathrm{core}(\widetilde{\Phi}_{\mathcal{C}})$
           (Theorem~\ref{thm:fuzzy-kvl}).
  \EndFor
  \For{each node $v_i\in V$}
    \State Apply $\mathcal{T}_{\mathrm{Back}}$: run Viterbi backward algorithm
           from current mode of $\tilde{\mu}_i$; compute
           $\dcat(\Back_i^{\mathrm{current}},\Back_i^H)$; restrict
           $\tilde{\mu}_i$ to concentrations consistent with backward
           trajectory within threshold.
  \EndFor
\Until{$d(\widetilde{\mathbf{X}},\widetilde{\mathbf{X}}^{\mathrm{prev}})<\varepsilon$}
\State \Return $\widetilde{\mathbf{X}}^*\leftarrow\widetilde{\mathbf{X}}$
\end{algorithmic}
\end{algorithm}

\subsection{Convergence Theorem}
\label{sec:alg-convergence}

\begin{theorem}[Convergence of Trajectory Completion]
\label{thm:convergence}
The operator $\mathcal{T}=\mathcal{T}_{\mathrm{Back}}\circ
\mathcal{T}_{\mathrm{KVL}}\circ\mathcal{T}_{\mathrm{KCL}}$ is a contraction
on $(\widetilde{\mathcal{X}},d)$ under the Hausdorff product metric.  By the
Banach fixed-point theorem~\cite{Banach1922}, iteration of $\mathcal{T}$
from any initial state converges to a unique fixed point
$\widetilde{\mathbf{X}}^*$.
\end{theorem}

\begin{proof}
We show each sub-operator is a strict contraction.

\emph{$\mathcal{T}_{\mathrm{KCL}}$ is a contraction.}
At node $v_i$, the KCL constraint relates $[C_i]$ linearly to the
concentrations of $v_i$'s neighbours.  For two fuzzy states
$\widetilde{\mathbf{X}}$ and $\widetilde{\mathbf{Y}}$, the resulting
KCL-consistent interval widths satisfy
\begin{align*}
  &d_H\bigl([\tilde{\mu}_i^{\mathrm{KCL}}(\mathbf{X})]^\alpha,
            [\tilde{\mu}_i^{\mathrm{KCL}}(\mathbf{Y})]^\alpha\bigr)\\
  &\quad\leq \sum_{j\in\mathcal{N}(i)}\frac{k_{ji}}{\sum_l k_{li}}\,
    d_H\bigl([\tilde{\mu}_j(\mathbf{X})]^\alpha,[\tilde{\mu}_j(\mathbf{Y})]^\alpha\bigr),
\end{align*}
where $\mathcal{N}(i)$ is the in-neighbour set of $v_i$ and the coefficients
$k_{ji}/(\sum_l k_{li})$ are positive and sum to at most 1.  Intersection with
the prior $\tilde{\mu}_i$ can only reduce the Hausdorff distance.  Hence
$\mathcal{T}_{\mathrm{KCL}}$ has contraction constant $c_1\leq 1$; it is
strictly $<1$ whenever at least one node is constrained by an observation.

\emph{$\mathcal{T}_{\mathrm{KVL}}$ is a contraction.}
The KVL constraint at each cycle imposes a multiplicative relation on
concentrations of cycle nodes.  The resulting restriction of the fuzzy
intervals narrows them toward the thermodynamic constraint manifold, reducing
the Hausdorff distance by a factor $c_2<1$ that depends on the cycle length
and the magnitude of $\Delta\mu^\circ$ values.

\emph{$\mathcal{T}_{\mathrm{Back}}$ is a contraction.}
The Viterbi backward trajectory from the current mode of $\tilde{\mu}_i$
identifies a set of ``backward-consistent'' concentrations.  Restricting
$\tilde{\mu}_i$ to this set reduces the support width monotonically.  The
Hausdorff distance between the restricted states satisfies
$d_H(\mathcal{T}_{\mathrm{Back}}(\tilde{\mu}),
\mathcal{T}_{\mathrm{Back}}(\tilde{\nu}))\leq c_3 d_H(\tilde{\mu},\tilde{\nu})$
with $c_3<1$ whenever the backward trajectory imposes a non-trivial constraint
(i.e., some concentrations are excluded).

\emph{Composition:} The composition of operators with contraction constants
$c_1,c_2,c_3$ gives overall constant $c\leq c_1 c_2 c_3<1$.  The space
$(\widetilde{\mathcal{X}},d)$ is complete because it is a finite product of
complete metric spaces (the Hausdorff metric on compact fuzzy sets is
complete)~\cite{Hausdorff1914}.  By the Banach fixed-point
theorem~\cite{Banach1922}, there exists a unique fixed point
$\widetilde{\mathbf{X}}^*$, and the iteration converges to it at the geometric
rate $c^n\,d(\widetilde{\mathbf{X}}_0,\widetilde{\mathbf{X}}^*)$.  $\square$
\end{proof}

\begin{corollary}[Crisp convergence under sufficient observations]
\label{cor:crisp}
If the observation set $\mathcal{O}$ is sufficient to uniquely determine the
network state (Definition~\ref{def:observability} below), then the fixed point
$\widetilde{\mathbf{X}}^*$ is a crisp state: $\tilde{\mu}_i^*$ has a
single-point support for all $i$.
\end{corollary}

\begin{proof}
Under sufficient observations, the three constraint operators jointly reduce
each fuzzy interval width to zero at convergence.  KCL propagates crisp
observations through the network; KVL eliminates thermodynamically inconsistent
assignments; the backward trajectory constraint eliminates kinetically
inconsistent assignments.  The fixed-point interval widths
$\overline{c}_i^{*,\alpha}-\underline{c}_i^{*,\alpha}\to 0$ for all $\alpha$
and all $i$, yielding crisp concentrations.  $\square$
\end{proof}

\subsection{Computational Complexity}
\label{sec:alg-complexity}

\begin{proposition}[Complexity of one algorithm iteration]
\label{prop:complexity}
For a network with $N$ species nodes, $R$ reaction edges, maximum vertex
degree $\Delta$, and backward trajectory depth $L$, one iteration of
Algorithm~\ref{alg:tc} requires:
\begin{itemize}[noitemsep]
  \item $O(N\Delta)$ for fuzzy KCL propagation;
  \item $O((R-N+1)\Delta)$ for fuzzy KVL (one pass per independent cycle);
  \item $O(N\cdot R\cdot L)$ for Viterbi backward trajectories.
\end{itemize}
Total: $O(N\cdot R\cdot L)$.  Convergence in $O(\log(1/\varepsilon)/\log(1/c))$
iterations gives total complexity $O(N\cdot R\cdot L\cdot\log(1/\varepsilon))$.
\end{proposition}

\begin{proof}
KCL at each node: $O(\Delta)$ fuzzy arithmetic operations for $O(\Delta)$
neighbours; $N$ nodes $\Rightarrow O(N\Delta)$.  KVL: $R-N+1$ independent
cycles by the circuit rank formula; each cycle traversal is $O(\Delta)$;
total $O((R-N+1)\Delta)\leq O(R\Delta)$.  Viterbi: at each of $N$ nodes, a
backward pass of depth $L$ with branching factor $\leq R$ requires $O(R\cdot L)$
per node, giving $O(N\cdot R\cdot L)$ total.  The geometric convergence bound
follows from Theorem~\ref{thm:convergence}.  $\square$
\end{proof}

\begin{remark}
For genome-scale metabolic models ($N,R\sim 10^3$--$10^4$; $\Delta\sim 10$;
$L\sim 20$--$100$), each iteration involves $\sim 10^8$--$10^{10}$ elementary
fuzzy arithmetic operations, feasible on modern hardware within minutes.  KCL
updates at different nodes can be parallelised (they are independent given
fixed neighbour states, analogous to Gauss-Seidel iteration), and Viterbi
backward passes for different nodes are completely independent, enabling
straightforward GPU parallelisation.
\end{remark}

%%=============================================================================
\section{Reactions as Time Generators}
\label{sec:time}
%%=============================================================================

\subsection{The Gillespie Framework}
\label{sec:time-gillespie}

The standard view of time in cellular dynamics treats it as an independent
real-valued parameter.  We propose instead that time is constituted by the
reaction firing sequence: there is no independent time axis; time IS the
ordered list of reaction events with associated waiting intervals.  This view
is already fully formalised by the Gillespie stochastic simulation
algorithm~\cite{Gillespie1977,Gillespie1976}.

\begin{definition}[Gillespie Stochastic Simulation]
\label{def:gillespie}
Given network state $\mathbf{n}$ and propensities $\{a_r(\mathbf{n})\}_{r=1}^R$,
the total propensity is $a_0=\sum_r a_r$.  The next reaction fires after
waiting time
\begin{equation}
  \tau = \frac{1}{a_0}\ln\frac{1}{u_1},\quad u_1\sim\mathrm{Unif}(0,1),
  \label{eq:gillespie-tau}
\end{equation}
and the firing reaction is
$r^*=\min\!\bigl\{r:\sum_{s=1}^r a_s\geq u_2\,a_0\bigr\}$,
$u_2\sim\mathrm{Unif}(0,1)$, independent of $u_1$~\cite{Gillespie1977}.
\end{definition}

\begin{proposition}[Reaction sequence is the complete causal record]
\label{prop:causal-record}
The complete stochastic trajectory of the reaction network is faithfully encoded
by the sequence of pairs $(R_k,\tau_k)_{k\geq 1}$, where $R_k\in\{1,\ldots,R\}$
is the identity of the $k$-th reaction event and $\tau_k>0$ is the preceding
waiting time.  The total elapsed time after $n$ events is
$T_n=\sum_{k=1}^n\tau_k$.
\end{proposition}

\begin{proof}
Given the initial state $\mathbf{n}_0$ and the sequence $(R_k,\tau_k)_{k\geq 1}$,
the state $\mathbf{n}(t)$ at any time $t$ is recovered by applying the
stoichiometric changes of reactions $R_1,R_2,\ldots$ in order until the first
event that occurs after $t$.  The sequence therefore encodes the full CTMC
trajectory; no additional information is needed.  $\square$
\end{proof}

\subsection{The Partition Lag}
\label{sec:time-lag}

Every elementary reaction requires a minimum time to execute: reactants must
encounter each other, reach the transition state, and reorganise into products.
We call this the \emph{partition lag}.

\begin{definition}[Partition Lag]
\label{def:partition-lag}
The \emph{partition lag} of reaction edge $e_{ij}$ is
\begin{equation}
  \taup_{ij} = \frac{\hbar}{\Delta E_{ij}} + \tau_{\mathrm{reorg},ij},
  \label{eq:partition-lag}
\end{equation}
where:
\begin{itemize}[noitemsep]
  \item $\hbar/\Delta E_{ij}$ is the quantum mechanical minimum transition time
        from the energy-time uncertainty relation
        $\Delta E\cdot\Delta t\geq\hbar/2$~\cite{Heisenberg1927}, with
        $\Delta E_{ij}$ the activation barrier;
  \item $\tau_{\mathrm{reorg},ij}$ is the network reorganisation time: the
        time required for the network to accommodate the state change triggered
        by the reaction (allosteric adjustments, diffusion of products,
        conformational relaxation).
\end{itemize}
\end{definition}

\begin{proposition}[Partition lag provides a strictly positive lower bound on inter-reaction intervals]
\label{prop:lag-minimum}
The Gillespie inter-event interval $\tau$ satisfies
$\tau\geq\taup_{\min}=\min_{(i,j)\in E}\taup_{ij}>0$.
\end{proposition}

\begin{proof}
The quantum bound $\hbar/\Delta E_{ij}>0$ for any finite $\Delta E_{ij}$.
Even for a formally barrierless reaction, $\tau_{\mathrm{reorg},ij}>0$ because
product release and downstream diffusion are non-instantaneous physical
processes.  The minimum over all edges is therefore strictly positive.
$\square$
\end{proof}

\begin{remark}
The partition lag is the ``clock tick'' of the cellular circuit.  Different
reaction classes operate on vastly different timescales:
\begin{itemize}[noitemsep]
  \item Electron transfer (Marcus theory~\cite{Marcus1956}):
        $\taup\sim 10^{-12}$--$10^{-9}$\,s.
  \item Proton transfer reactions:
        $\taup\sim 10^{-11}$--$10^{-8}$\,s.
  \item Enzymatic catalysis~\cite{Wolfenden2001,MichaelisMenten1913}:
        $\taup\sim 10^{-6}$--$10^{-3}$\,s.
  \item Transcription and translation (gene expression):
        $\taup\sim 10^2$--$10^4$\,s.
\end{itemize}
The resulting hierarchy spans twelve orders of magnitude, providing the
physical basis for multi-timescale decomposition.
\end{remark}

\subsection{Heterogeneous Cellular Time}
\label{sec:time-hetero}

\begin{definition}[Cellular Time Hierarchy]
\label{def:time-hierarchy}
Partition the reaction set $E$ into classes by partition lag:
\begin{align*}
  E_{\mathrm{fast}}   &= \{e_{ij}:\taup_{ij}<10^{-6}\,\text{s}\},\\
  E_{\mathrm{medium}} &= \{e_{ij}:10^{-6}\leq\taup_{ij}<1\,\text{s}\},\\
  E_{\mathrm{slow}}   &= \{e_{ij}:\taup_{ij}\geq 1\,\text{s}\}.
\end{align*}
\end{definition}

\begin{theorem}[Multi-timescale reduction]
\label{thm:timescale}
When $\max_{e\in E_{\mathrm{fast}}}\taup_e\ll\min_{e\in E_{\mathrm{medium}}}\taup_e\ll
\min_{e\in E_{\mathrm{slow}}}\taup_e$, the fast reactions reach quasi-steady
state before any medium-timescale reaction fires, and medium reactions reach
quasi-steady state before any slow reaction fires.  The dynamics on each
timescale can be analysed by a reduced system with QSS values from faster
timescales as fixed parameters~\cite{Murray2002,HeinrichSchuster1996}.
\end{theorem}

\begin{proof}
Standard singular perturbation / quasi-steady-state (QSS) reduction~\cite{Murray2002}.
Let $\varepsilon=\max(\taup_{\mathrm{fast}})/\min(\taup_{\mathrm{medium}})\to 0$.
Writing the fast variables as $\mathbf{y}$ and the slow variables as
$\mathbf{z}$, the system decomposes as $\varepsilon\dot{\mathbf{y}}=
f(\mathbf{y},\mathbf{z})$, $\dot{\mathbf{z}}=g(\mathbf{y},\mathbf{z})$.
On the medium timescale, $\varepsilon\to 0$ gives the algebraic QSS constraint
$f(\mathbf{y}^{\mathrm{QSS}},\mathbf{z})=0$, eliminating $\mathbf{y}$.
Iterating this argument across the hierarchy gives the multi-level reduction.
$\square$
\end{proof}

Figure~\ref{fig:timescales} illustrates the cellular time hierarchy.



%%=============================================================================
\section{Categorical State Propagation}
\label{sec:propagation}
%%=============================================================================

\subsection{Signal Velocity vs.\ Drift Velocity}
\label{sec:prop-velocity}

A central question for the circuit model is whether the network is
informationally well-mixed: does categorical state information propagate
through the network faster than molecular drift?  If $v_s\gg v_d$, the global
circuit description is appropriate; if $v_s\sim v_d$, spatial effects dominate.

\begin{definition}[Signal Velocity and Drift Velocity]
\label{def:velocities}
In a coupled molecular network:
\begin{itemize}[noitemsep]
  \item The \emph{signal velocity} $v_s$ is the speed at which a perturbation
        in the state of one node propagates to adjacent nodes.  For
        mechanically coupled networks (allosteric signalling), $v_s\sim
        \sqrt{\kappa/m_{\mathrm{eff}}}$, where $\kappa$ is the effective
        coupling stiffness and $m_{\mathrm{eff}}$ is the effective molecular
        mass.  For collision-mediated propagation,
        $v_s\sim k_{\mathrm{on}}[C]\cdot d$, where $d$ is the molecular
        encounter distance and $k_{\mathrm{on}}$ is the bimolecular
        association rate.
  \item The \emph{drift velocity} $v_d$ is the characteristic Fickian
        transport speed: $v_d\sim D/L$, where $D$ is the cytoplasmic
        diffusion coefficient and $L$ is the cell length.
\end{itemize}
\end{definition}

\begin{theorem}[Network is informationally well-mixed]
\label{thm:well-mixed}
For a typical intracellular molecular network, the signal-to-drift velocity
ratio satisfies
\begin{equation}
  \frac{v_s}{v_d} \;\sim\; \frac{L\sqrt{\kappa/m_{\mathrm{eff}}}}{D}
    \;\approx\; 10^3\text{--}10^6,
  \label{eq:velocity-ratio}
\end{equation}
so categorical state information propagates three to six orders of magnitude
faster than molecular drift.  The network is informationally well-mixed on
any timescale longer than the partition lag $\taup_{\min}$.
\end{theorem}

\begin{proof}
\emph{Allosteric signal velocity:}  Conformational changes in proteins
propagate at the speed of sound in condensed matter.  Taking Young's modulus
$E_Y\approx 10^9$\,Pa and density $\rho\approx 10^3$\,kg\,m$^{-3}$,
$v_s^{\mathrm{allosteric}}\approx\sqrt{E_Y/\rho}\approx 10^3$\,m\,s$^{-1}$.

\emph{Collision-mediated signal velocity:}  At a typical intracellular
metabolite concentration $[C]\approx 1\,\mu$M with association rate
$k_{\mathrm{on}}\approx 10^9$\,M$^{-1}$\,s$^{-1}$ (diffusion-limited), the
per-molecule collision rate is $k_{\mathrm{on}}[C]\approx 10^3$\,s$^{-1}$.
Each collision propagates a state change over an encounter distance
$d\approx 1$\,nm, giving $v_s^{\mathrm{collision}}\approx
10^3\,\mathrm{s}^{-1}\times 10^{-9}\,\mathrm{m}=10^{-6}$\,m\,s$^{-1}=
1\,\mu$m\,s$^{-1}$.

\emph{Drift velocity:}  For a metabolite in cytoplasm,
$D\approx 10^{-10}$\,m$^2$\,s$^{-1}$; with $L\approx 10^{-5}$\,m
(typical mammalian cell), $v_d=D/L\approx 10^{-5}$\,m\,s$^{-1}=10\,\mu$m\,s$^{-1}$.
For a protein, $D\approx 10^{-11}$\,m$^2$\,s$^{-1}$,
$v_d\approx 10^{-6}$\,m\,s$^{-1}=1\,\mu$m\,s$^{-1}$.

\emph{Ratios:}  Allosteric vs.\ protein drift:
$v_s^{\mathrm{allosteric}}/v_d^{\mathrm{protein}}\approx 10^3/10^{-6}=10^9$;
collision-mediated vs.\ metabolite drift:
$v_s^{\mathrm{collision}}/v_d^{\mathrm{metabolite}}\approx
10^{-6}/10^{-5}=10^{-1}$, but at higher concentrations
($[C]\approx 1$\,mM: $k_{\mathrm{on}}[C]\approx 10^6$\,s$^{-1}$, so
$v_s\approx 10^{-3}$\,m\,s$^{-1}$) the ratio becomes $10^2$.
Across the range of biological conditions,~\eqref{eq:velocity-ratio}
gives $10^3$--$10^6$.  $\square$
\end{proof}

\begin{corollary}[Global network structure over local concentration]
\label{cor:global}
Because $v_s\gg v_d$, categorical state information equilibrates across the
entire network before any significant net change in average concentration
occurs.  Therefore: (i) the effective size of the cell in information space is
much smaller than in physical space; (ii) global network topology (captured by
the circuit model) is the appropriate level of description; (iii) local
concentration gradients are secondary to circuit topology for determining node
state.
\end{corollary}

\begin{proof}
The time for a signal to traverse the network is $t_s\sim L/v_s$.  The time
for a concentration gradient of magnitude $\Delta C$ to build up by Fickian
diffusion is $t_d\sim L^2/(D)=L/v_d$.  Since $v_s\gg v_d$,
$t_s\ll t_d$: the signal arrives before any significant drift can occur,
establishing informational equilibration.  $\square$
\end{proof}

\begin{figure*}[!htbp]
    \centering
    \includegraphics[width=\textwidth]{figures/panel_exp5.png}
    \caption{Categorical Signal Velocity vs.\ Material Drift (Theorem~8.1).
    (a)~Signal propagation velocity $v_{\mathrm{sig}}$ for nine enzymatic
    mechanisms and hydrogen-bond proton transport on a logarithmic scale;
    enzymatic signals range from $10^{-3}$ to $2 \times 10^{-2}\;\mathrm{m\,s^{-1}}$,
    while H-bond categorical propagation reaches
    ${\sim}7 \times 10^{13}\;\mathrm{m\,s^{-1}}$.
    (b)~$\log_{10}$ of the signal-to-drift velocity ratio per mechanism;
    enzymatic channels exceed drift by $10^2$--$10^{3.3}\times$, while the
    H-bond network achieves a ratio of ${\sim}10^{17.5}$, consistent with
    Newton's-cradle categorical state propagation.
    (c)~Three-dimensional scatter of enzymatic mechanisms in
    ($k_{\mathrm{cat}}$, $v_{\mathrm{sig}}$, ratio) space, showing the linear
    scaling $v_{\mathrm{sig}} = k_{\mathrm{cat}} \cdot a_{\mathrm{mol}}$ and
    the proportional growth of the velocity ratio with turnover number.
    (d)~Log--log plot of $v_{\mathrm{sig}}$ vs.\ $v_{\mathrm{drift}}$ for all
    mechanisms; the grey diagonal marks $v_{\mathrm{sig}} = v_{\mathrm{drift}}$;
    all points lie above this line by two or more decades, confirming that
    categorical state propagation universally dominates material transport
    ($v_{\mathrm{sig}} / v_{\mathrm{drift}} \geq 100$).}
    \label{fig:panel_exp5}
\end{figure*}

\subsection{Reversibility from Categorical Flow Direction}
\label{sec:prop-reversibility}

\begin{theorem}[Reaction direction from categorical depth gradient]
\label{thm:direction}
A reaction edge $e_{ij}$ carries net positive flux ($J_{ij}>0$, forward
direction) if and only if
\begin{equation}
  \Hcat_j(\sigma_j) - \Hcat_i(\sigma_i) \;>\;
    \frac{c_j - c_i}{\kB T\ln 2},
  \label{eq:direction}
\end{equation}
i.e., the reaction is thermodynamically spontaneous.  At thermodynamic
equilibrium ($J_{ij}=0$), categorical state flows bidirectionally with equal
magnitude in both directions---detailed balance---corresponding to zero net
information flux but continuous molecular turnover.
\end{theorem}

\begin{proof}
By Corollary~\ref{cor:ohm}, near equilibrium
$J_{ij}\approx\Gcond_{ij}(\phi_i-\phi_j^{\mathrm{eff}})$.  Substituting
$\phi_i=-\kB T\ln 2\cdot\Hcat_i+c_i$ (Theorem~\ref{thm:mu-catdepth}):
\[
  J_{ij}>0 \Leftrightarrow \phi_i > \phi_j^{\mathrm{eff}}
  \Leftrightarrow -\kB T\ln 2\,\Hcat_i+c_i > -\kB T\ln 2\,\Hcat_j+c_j,
\]
giving~\eqref{eq:direction} after rearrangement.  At equilibrium, detailed
balance~\cite{Lewis1925} requires $k_{ij}[C_i]^{\mathrm{eq}}=
k_{ji}[C_j]^{\mathrm{eq}}$, so $J_{ij}=0$ and forward/reverse fluxes are
equal; categorical state flows in both directions at equal rate.  $\square$
\end{proof}

\begin{remark}
Theorem~\ref{thm:direction} provides the information-theoretic explanation for
reaction reversibility.  A reaction is reversible not because the molecular
mechanism changes but because the direction of categorical state flow---the
direction in which information propagates through the circuit---can reverse
when the free energy gradient reverses.  The reaction topology (the edge in
$G$) is invariant; only the direction of net current changes.
\end{remark}

\subsection{Protein Isoforms as Degenerate Topologies}
\label{sec:prop-isoforms}

\begin{definition}[Isoform-Degenerate Edges]
\label{def:isoforms}
Two edges $e_{ij}^{(1)}$ and $e_{ij}^{(2)}$ in $G$ are
\emph{isoform-degenerate} if they connect the same substrate and product nodes
$(v_i,v_j)$ via distinct intermediate complex nodes $v_c^{(1)},v_c^{(2)}$
(corresponding to different isoforms of the enzyme catalysing the reaction).
They implement the same categorical state transition $\sigma_i\to\sigma_j$
via different molecular mechanisms.
\end{definition}

\begin{proposition}[Isoforms act as parallel conductances]
\label{prop:isoforms}
For $m$ isoform-degenerate edges $e_{ij}^{(1)},\ldots,e_{ij}^{(m)}$, the
effective conductance is
\begin{equation}
  \Gcond_{ij}^{\mathrm{eff}} = \sum_{k=1}^m \Gcond_{ij}^{(k)},
  \label{eq:isoform-cond}
\end{equation}
in exact analogy with parallel resistors.
\end{proposition}

\begin{proof}
The total net flux from $v_i$ to $v_j$ via all isoforms is additive (the
isoforms do not interact in this limit):
$J_{ij}^{\mathrm{total}}=\sum_k J_{ij}^{(k)}=\sum_k\Gcond_{ij}^{(k)}\Delta\phi_{ij}
=(\sum_k\Gcond_{ij}^{(k)})\Delta\phi_{ij}=\Gcond_{ij}^{\mathrm{eff}}\Delta\phi_{ij}$,
since all isoforms respond to the same potential difference $\Delta\phi_{ij}$.
$\square$
\end{proof}

\begin{corollary}[Isoform redundancy provides network robustness]
\label{cor:isoform-robust}
Even if one isoform is inactivated ($\Gcond_{ij}^{(k)}\to 0$), the effective
conductance remains positive as long as at least one isoform is active.  The
categorical state transition continues; only its rate changes.  This is the
molecular basis for the well-known robustness of metabolic networks to enzyme
knockouts~\cite{BarabasiOltvai2004,Alon2007book}.
\end{corollary}

%%=============================================================================
\section{Applications}
\label{sec:applications}
%%=============================================================================

\subsection{Cell-State Instantiation from Partial Observations}
\label{sec:app-instantiation}

\emph{Problem setting.}  Given a network $G$, standard potentials
$\{\mu_i^\circ\}$, and partial observations
$\mathcal{O}=\{[C_{i_k}]\pm\epsilon_k\}_{k=1}^M$ from an omics experiment
(proteomics, metabolomics, transcriptomics), determine the complete molecular
state $\{[C_i]\}_{i=1}^N$.

\emph{Solution.}  Apply Algorithm~\ref{alg:tc}.  Observed nodes initialise
as narrow trapezoidal fuzzy states.  Unobserved nodes initialise to the
maximum-entropy prior: $\tilde{\mu}_i(c)\equiv 1$ on the physiological range
$[c_{\min}^i,c_{\max}^i]$ from databases (BRENDA, BiGG, HMDB).  KCL
propagation tightens unobserved nodes using observed nodes as anchors.  KVL
eliminates thermodynamically inconsistent assignments.  The backward trajectory
constraint adds kinetic consistency.  Theorem~\ref{thm:convergence} guarantees
convergence to $\widetilde{\mathbf{X}}^*$.

\begin{definition}[Observability]
\label{def:observability}
A network $G$ is \emph{observable} from $\mathcal{O}$ if
Algorithm~\ref{alg:tc} converges to a crisp state (Corollary~\ref{cor:crisp}).
\end{definition}

\begin{proposition}[Sufficient condition for observability]
\label{prop:observability}
If $\mathcal{O}$ forms a dominating set of $G$ (every unobserved node has at
least one observed neighbour), then $G$ is observable from $\mathcal{O}$
under the three constraint classes.
\end{proposition}

\begin{proof}
If every unobserved node $v_i$ has an observed neighbour, KCL propagation at
$v_i$ (Theorem~\ref{thm:fuzzy-kcl}) constrains $[C_i]$ to a unique interval
given the observed neighbour concentrations and the rate constants.  Subsequent
KVL enforcement (Theorem~\ref{thm:fuzzy-kvl}) eliminates thermodynamically
inconsistent candidates, and the backward trajectory constraint (Definition~\ref{def:operators})
further restricts to kinetically consistent candidates.  Under the combined
constraints, the interval width converges to zero, giving a crisp value.
Iterating over all unobserved nodes gives the complete crisp state.  $\square$
\end{proof}

\begin{example}[Glycolytic state instantiation]
\label{ex:glycolysis}
Consider the glycolytic pathway ($N=10$ primary metabolites,
$R=10$ enzymatic reactions).  Suppose glucose, ATP, ADP, and NADH are
measured ($M=4$).  Every glycolytic intermediate (fructose-6-phosphate,
fructose-1,6-bisphosphate, etc.) has at least one neighbour among the observed
species via the pathway graph.  Algorithm~\ref{alg:tc} therefore converges to
the complete glycolytic state, including unmeasured intermediates, consistent
with all mass-balance, thermodynamic, and kinetic constraints.
\end{example}

\subsection{Disease Detection from Trajectory Analysis}
\label{sec:app-disease}

\emph{Problem setting.}  Given partial observations from a cell of unknown
health status, determine whether the cell is healthy or diseased.

\emph{Solution.}  First run Algorithm~\ref{alg:tc} to obtain
$\widetilde{\mathbf{X}}^*$.  Then compute the backward trajectory
$\Back_i^{\mathrm{obs}}$ of each node from the fixed-point state and compare
against the healthy reference $\Back_i^H$.

\begin{definition}[Trajectory Escape]
\label{def:traj-escape}
Node $v_i$ exhibits \emph{trajectory escape} if
\begin{equation}
  \dcat\bigl(\Back_i^{\mathrm{obs}},\,\Back_i^H\bigr) > \theta_i,
  \label{eq:escape}
\end{equation}
where $\theta_i$ is a per-node threshold set at three standard deviations above
the healthy baseline distribution of $\dcat$.
\end{definition}

\begin{theorem}[Disease detection by trajectory escape]
\label{thm:disease-detect}
Under the assumptions of Theorem~\ref{thm:health-address}, a cell is diseased
if and only if at least one node exhibits trajectory escape, and healthy if no
node does.
\end{theorem}

\begin{proof}
By Theorem~\ref{thm:health-address}, a diseased cell has at least one node
whose backward trajectory converges to $\mathcal{A}_D\neq\mathcal{A}_H$.  The
categorical distance $\dcat(\Back_i^{\mathrm{obs}},\Back_i^H)$ quantifies the
divergence.  When $\dcat>\theta_i$ (set to capture $>99.7\%$ of healthy
variability under Gaussian assumptions), it signals escape from $\mathcal{B}_H$.
Conversely, $\dcat<\theta_i$ for all nodes implies that all backward
trajectories are consistent with $\mathcal{A}_H$.  $\square$
\end{proof}

\begin{remark}
This criterion does not require measuring every molecular species.  Even sparse
observations, when propagated through the trajectory completion algorithm,
yield a complete state estimate from which backward trajectories can be
computed.  A biomarker is formally defined as a species with high
\emph{trajectory sensitivity} $\Xi_i=\partial\,\dcat/\partial\|\delta\mathbf{x}_D\|$:
small perturbations to its concentration strongly shift its backward trajectory,
enabling early detection with few measurements.
\end{remark}

\subsection{Rational Drug Design as Edge-Conductance Modification}
\label{sec:app-drugs}

In the circuit model, a drug targeting enzyme $e_{ij}$ modifies the
conductance $\Gcond_{ij}$.

\begin{definition}[Drug Action as Conductance Perturbation]
\label{def:drug}
A drug $D$ targeting edge $e_{ij}$ modifies
$\Gcond_{ij}\mapsto\Gcond_{ij}(1-\eta_D)$, where $\eta_D\in[0,1]$ is the
fractional inhibition (activation if $\eta_D<0$).  For a competitive inhibitor
with dissociation constant $K_I$, the Michaelis-Menten conductance becomes
\begin{equation}
  \Gcond_{SP}^{\mathrm{inh}} = \frac{k_{\mathrm{cat}}[E]_T}
    {RT(K_m(1+[D]/K_I)+[S])}.
  \label{eq:mm-inhib}
\end{equation}
\end{definition}

\begin{definition}[Minimal Drug Design Criterion]
\label{def:drug-design}
An optimal drug restores health by modifying conductances so that trajectory
completion yields a fixed point $\widetilde{\mathbf{X}}^{*,\mathrm{treated}}$
with $\dcat(\Back_i^{\mathrm{obs,treated}},\Back_i^H)<\theta_i$ for all $i$,
while minimising the total conductance perturbation (minimal pharmacological
load):
\begin{equation}
  \min_{\{\eta_{ij}\}}\sum_{ij}|\eta_{ij}|
  \quad\text{s.t.}\quad
  \dcat(\Back_i^{\mathrm{obs,treated}},\Back_i^H)<\theta_i\;\forall\,i,
  \quad\eta_{ij}\in[0,1].
  \label{eq:drug-opt}
\end{equation}
\end{definition}

\begin{proposition}[Drug design as a sparse inverse conductance problem]
\label{prop:drug-inverse}
Under first-order linearisation of the backward-trajectory distance in the
$\{\eta_{ij}\}$, the optimisation~\eqref{eq:drug-opt} is a linear programme
with an $\ell_1$ objective, solvable in polynomial time.  The $\ell_1$
objective promotes sparsity, identifying the smallest set of reactions that
need to be targeted.
\end{proposition}

\begin{proof}
Linearising $\dcat(\Back_i^{\mathrm{obs,treated}},\Back_i^H)$ about the
diseased state $\{\eta_{ij}=0\}$ gives a linear function of $\{\eta_{ij}\}$
plus higher-order terms.  In the linear approximation, the feasibility
constraints become $\sum_{ij}A_{i,ij}\eta_{ij}\geq b_i$ (where $A_{i,ij}$ and
$b_i$ encode the linearised trajectory-distance constraints), which is a
system of linear inequalities.  The objective $\sum_{ij}|\eta_{ij}|$ is the
$\ell_1$ norm, cast as a standard LP by introducing variables $s_{ij}\geq
|\eta_{ij}|$: minimise $\sum_{ij}s_{ij}$ subject to
$s_{ij}\geq\pm\eta_{ij}$ and the feasibility constraints.
This LP is solvable in polynomial time~\cite{Pearl1988}, and the $\ell_1$
objective produces a sparse solution by the standard compressed-sensing
argument.  $\square$
\end{proof}

\begin{figure*}[!htbp]
    \centering
    \includegraphics[width=\textwidth]{figures/panel_exp6.png}
    \caption{Disease Detection via Backward-Trajectory Divergence
    (Theorem~9.2).
    (a)~Steady-state concentrations (log scale) for healthy, hexokinase-impaired
    (HK$_{10\%}$), and phosphofructokinase-impaired (PFK$_{10\%}$) circuits
    across all 10 glycolytic nodes; HK impairment causes $43\times$ GLC
    accumulation while downstream intermediates collapse by
    $10$--$100\times$; PFK impairment produces $128\times$ G6P/F6P buildup
    upstream of the block.
    (b)~Log--log scatter of disease vs.\ healthy steady-state concentrations
    for HK impairment; deviation from the identity line (grey) identifies
    GLC as the escaped node, with maximal relative deviation $4{,}233\%$.
    (c)~Three-dimensional projection of
    $(\log_{10} c_{\mathrm{healthy}},\; \log_{10} c_{\mathrm{HK}},\;
    \log_{10} c_{\mathrm{PFK}})$ showing the distinct displacement patterns
    of the two disease conditions in metabolic state space; healthy-proximal
    nodes cluster near the origin while escaped nodes scatter to the
    periphery.
    (d)~Circuit coherence $\eta$ on a logarithmic scale: healthy circuit
    maintains $\eta = 1.0$ (full phase-lock); HK impairment reduces coherence
    to $\eta \approx 0$ and PFK impairment to $\eta \sim 10^{-49}$,
    confirming that both disease states correspond to complete loss of
    backward-trajectory closure (basin escape).}
    \label{fig:panel_exp6}
\end{figure*}

%%=============================================================================
\section{Discussion}
\label{sec:discussion}
%%=============================================================================

\subsection{Relationship to Existing Approaches}
\label{sec:disc-existing}

\emph{Flux Balance Analysis (FBA).}
FBA~\cite{Orth2010,Palsson2006,Price2004} maximises an objective flux subject
to the steady-state constraint $\mathbf{S}\mathbf{v}=0$ and flux bounds.  The
biochemical KCL (Theorem~\ref{thm:kcl}) is identical to the FBA steady-state
constraint.  The circuit model extends FBA by: (i) incorporating thermodynamic
KVL constraints, eliminating thermodynamically infeasible flux
distributions~\cite{Price2004}; (ii) providing a Bayesian/Viterbi inferential
framework rather than a purely optimisation-based one; (iii) accommodating
fuzzy inputs rather than requiring crisp bounds; (iv) inferring concentrations
(not only fluxes) via backward trajectory analysis.

\emph{Metabolic Control Analysis (MCA).}
MCA~\cite{KacserBurns1973,FellSauro1985,Fell1992,HeinrichSchuster1996}
characterises flux and concentration control coefficients.  In the circuit
model, the flux control coefficient $C^J_e=\partial\ln J/\partial\ln v_e$
is the logarithmic derivative of circuit current with respect to edge
conductance, exactly analogous to the sensitivity of a resistor network.
Proposition~\ref{prop:drug-inverse} generalises the MCA summation and
connectivity theorems to the inverse problem: given a desired flux change,
identify the minimal enzyme perturbation.

\emph{Kinetic modelling.}
Explicit kinetic ODE models~\cite{Murray2002,Klipp2005,Voit2000book} require
precise values of all rate constants and initial conditions---rarely available
for complex networks.  The circuit model with fuzzy states handles uncertain
parameters naturally, and trajectory completion infers states without requiring
the full kinetic parameter set.

\emph{Constraint-based reconstruction and analysis (COBRA).}
COBRA methods~\cite{Palsson2015,Price2004} use the stoichiometric matrix plus
flux bounds to constrain the feasible flux space.  The circuit model incorporates
all COBRA constraints (via KCL) and adds three additional constraint types (KVL,
backward trajectory, thermodynamic feasibility) that further restrict the
feasible state space and uniquely identify concentrations.

\emph{Network biology.}
The informational well-mixing result (Theorem~\ref{thm:well-mixed}) connects to
the small-world architecture of metabolic and protein interaction
networks~\cite{BarabasiOltvai2004,Wagner2001}: short average path lengths
ensure that categorical state information propagates quickly across the network.
Network motifs~\cite{Alon2007,Alon2007book} (feedforward loops, autoregulatory
feedback) correspond to specific subgraph topologies in the circuit model, each
with a characteristic backward trajectory pattern.

\subsection{The Emergence of Directed Reactions and Autocatalytic Closure}
\label{sec:disc-autopoiesis}

A deeper implication of the framework concerns the origin of directed reactions.
In classical kinetics, reaction direction is determined by the sign of $\Delta G$:
reactions run downhill.  But why does the free energy landscape have the shape
it does?

In the circuit model, $\phi_i=-\kB T\ln 2\cdot\Hcat_i+c_i$ and $\Hcat_i$ is
determined by $P(\sigma_i)$, the stationary probability under the network
dynamics.  The stationary distribution itself is determined by the network
topology and rate constants.  The free energy landscape and reaction directions
are therefore \emph{self-consistently determined} by the network's own topology:
the network shapes its own energy landscape.

This is the information-theoretic formulation of what Kauffman terms
\emph{autocatalytic closure}~\cite{Kauffman1993}: the network catalyses its
own maintenance.  Varela, Maturana, and Uribe~\cite{VarelaMaturanaUribe1974}
call this \emph{autopoiesis}: the network is simultaneously the product of and
the producer of its own components.  Rosen~\cite{Rosen1991} identifies this
self-referential property---that the formal cause of the network is internal to
the network itself---as the defining characteristic distinguishing living from
non-living systems.

In the circuit model, autocatalytic closure is expressed formally as follows:
the MAP backward trajectories of all nodes close on the network itself---every
backward path eventually returns to a node already in the network, rather than
escaping to an external attractor.  This is precisely the definition of the
healthy state in Theorem~\ref{thm:health-address}.  Disease is the breakdown
of autocatalytic closure: some backward trajectories escape, meaning that the
cell can no longer self-consistently maintain those nodes within their healthy
attractor basin.

A canonical example of energy transduction in biological circuits is the
chemiosmotic mechanism of Mitchell~\cite{Mitchell1961}: the proton gradient
across the inner mitochondrial membrane acts as a battery (EMF source in
Table~\ref{tab:analogy}), driving ATP synthesis.  In the circuit model, the
proton-motive force is the potential difference $\Delta\phi_{H^+}=
\Delta\psi - (RT/F)\Delta\mathrm{pH}$, exactly a node potential difference,
and the ATP synthase is a conductance element coupling the $H^+$ current to
the phosphorylation flux~\cite{NichollsFerguson2002}.

The Onsager reciprocal relations~\cite{Onsager1931} emerge naturally in the
near-equilibrium regime: the linearised conductance matrix $\Gcond_{ij}$ is
symmetric (since $\Gcond_{ij}$ and $\Gcond_{ji}$ are related by detailed
balance), satisfying $L_{ij}=L_{ji}$ without additional postulates.  Far from
equilibrium, the nonlinear Michaelis-Menten conductances break Onsager
symmetry and give rise to dissipative structures~\cite{PrigogineNicolis1977,
Prigogine1967}: sustained oscillations (as in glycolysis~\cite{EpsteinPojman1998},
the cell cycle~\cite{Tyson2003}, and the Lotka-Volterra models of population
ecology~\cite{Lotka1925,Volterra1926,Kolmogorov1936}), spatial patterns, and
deterministic chaos~\cite{Strogatz1994}.  In the circuit model, these
correspond to regimes where the fixed point of Algorithm~\ref{alg:tc} is not
a single concentration assignment but a fuzzy set supported on a limit cycle
or strange attractor.

Self-organisation of the reaction network~\cite{Eigen1971} into
functional modules---pathways, signalling cascades, gene regulatory
circuits---appears in the circuit model as the emergence of modular subgraphs
with strong internal conductances and weak inter-module conductances.  The
information flow is then predominantly intra-modular, with categorical addresses
encoding module membership.

\subsection{Limitations and Future Work}
\label{sec:disc-limitations}

\emph{Spatial heterogeneity.}
Theorem~\ref{thm:well-mixed} establishes informational well-mixing in a single
compartment.  For polarised cells, highly localised signalling platforms
(postsynaptic densities, ciliary compartments), or large cells (neurons with
metre-long axons), $v_s/v_d$ may approach unity and spatial effects become
important.  Extension to a spatial circuit---a continuum of circuit elements
distributed in three-dimensional space, connected by diffusion
resistors~\cite{Murray2002}---is needed.

\emph{Quantum effects.}
The partition lag includes a quantum uncertainty term.  For reactions involving
proton transfer or electron tunnelling~\cite{Marcus1956}, quantum coherence
effects may persist at biological temperatures, and the classical circuit model
would need augmentation with quantum circuit theory.

\emph{Parameter identifiability.}
The conductances $\Gcond_{ij}$ require knowledge of rate constants, which are
poorly known for most cellular reactions.  A full treatment would require
Bayesian inference over both states and parameters simultaneously~\cite{Pearl1988}.

\emph{Genome-scale validation.}
Systematic validation on genome-scale metabolic reconstructions (RECON3D,
iML1515)~\cite{Palsson2015} is needed to characterise practical convergence
rates and minimum observation sets for crisp state determination.

\emph{Dynamic disease trajectories.}
The backward trajectory framework treats a snapshot.  Many diseases are
dynamical processes: the cell transitions over time.  Extending
Theorem~\ref{thm:health-address} to time-varying networks would require a
non-autonomous backward trajectory, losing the time-invariance of
Theorem~\ref{thm:timeinv} and requiring additional temporal constraints.

\emph{Cell-cell communication.}
Cells in tissues communicate via signalling molecules, gap junctions, and
tunnelling nanotubes.  The single-cell circuit model would need extension to a
network-of-networks~\cite{BarabasiOltvai2004}, with inter-cell edges connecting
secreted products of one cell to the receptors of another.

%%=============================================================================
\section{Conclusion}
\label{sec:conclusion}
%%=============================================================================

We have presented a unified cellular circuit model in which the biochemical
reaction network of a living cell is formalised as an electrical circuit graph.
Every element of the analogy is derived from first principles: node potentials
from the chemical potential (Theorem~\ref{thm:mu-catdepth}), edge currents from
reaction fluxes (Definition~\ref{def:current}), edge conductances from
Eyring--Arrhenius and Michaelis-Menten kinetics
(Definitions~\ref{def:arrhenius-cond} and~\ref{def:mm-cond}), KCL from
stoichiometric mass balance (Theorem~\ref{thm:kcl}), and KVL from
thermodynamic cycle consistency (Theorem~\ref{thm:kvl}).  The chemical potential
at each node equals the information-theoretic categorical depth up to the
thermal energy scale $\kB T\ln 2$, establishing a rigorous equivalence between
thermodynamic and information-theoretic descriptions of cellular state.

Epistemic blindness (Proposition~\ref{prop:blind})---the cell's inability to
self-diagnose---motivates the shift from forward simulation to trajectory
completion.  The trajectory completion algorithm (Algorithm~\ref{alg:tc}) solves
the inverse problem of cellular state determination from partial observations by
iterating three constraint classes: fuzzy Kirchhoff laws
(Theorems~\ref{thm:fuzzy-kcl} and~\ref{thm:fuzzy-kvl}), MAP backward trajectory
inference via the Viterbi algorithm~\cite{Viterbi1967,Forney1973}, and
thermodynamic feasibility.  Convergence to a unique fixed point is guaranteed
by the Banach fixed-point theorem~\cite{Banach1922}
(Theorem~\ref{thm:convergence}).

The MAP backward trajectory of each node is time-invariant for autonomous
kinetics (Theorem~\ref{thm:timeinv}), making it a categorical address: a fixed
geometric identity of the node's state in the network's state space.  Disease
is detected when a node's backward trajectory escapes the healthy attractor
basin (Theorem~\ref{thm:health-address}), providing a principled diagnostic
criterion.  Drug design reduces to a convex $\ell_1$ optimisation problem:
find the minimal edge-conductance modification that restores escaped backward
trajectories to the healthy basin (Proposition~\ref{prop:drug-inverse}).

Categorical state information propagates $10^3$--$10^6\times$ faster than
molecular drift (Theorem~\ref{thm:well-mixed}), ensuring informational
well-mixing and validating the global circuit description.  Reaction
reversibility is explained as bidirectional categorical flow at equilibrium
(Theorem~\ref{thm:direction}), and protein isoforms are degenerate parallel
implementations of the same categorical state transition
(Proposition~\ref{prop:isoforms}).  Autocatalytic closure
(Kauffman~\cite{Kauffman1993}, Varela--Maturana~\cite{VarelaMaturanaUribe1974},
Rosen~\cite{Rosen1991}) corresponds to the condition that backward trajectories
of all nodes converge within the network, precisely the healthy-state criterion
of Theorem~\ref{thm:health-address}.

The framework unifies flux balance analysis, metabolic control analysis,
stochastic kinetics, and information theory within a single information-geometric
picture, while adding new capabilities: fuzzy state inference, backward
trajectory computation, time-invariant categorical addressing, and
trajectory-escape-based disease classification.  We anticipate that these
capabilities will find application in precision medicine (personalised biomarker
panels), synthetic biology (designing circuits that maintain specific attractor
basins), and fundamental systems biology (characterising the information geometry
of cellular state space).

%%=============================================================================
\bibliographystyle{plainnat}
\bibliography{references}
%%=============================================================================

\end{document}
